\documentclass[aps,prd,twocolumn,superscriptaddress,nofootinbib]{revtex4-2}

\usepackage[T1]{fontenc}
\usepackage[utf8]{inputenc}
\input{glyphtounicode}
\pdfgentounicode=1
\usepackage{lmodern}
\usepackage[ngerman]{babel}
\usepackage{amsmath}
\usepackage{mathtools}
\usepackage{amssymb}
\usepackage{amsthm}
\usepackage{booktabs}
\usepackage{array}
\usepackage{xcolor}
\usepackage{graphicx}
\usepackage{mathrsfs}

\emergencystretch=2em

\newtheorem{theorem}{Satz}[section]
\newtheorem{proposition}[theorem]{Proposition}
\newtheorem{lemma}[theorem]{Lemma}
\newtheorem{corollary}[theorem]{Korollar}
\newtheorem{conjecture}[theorem]{Vermutung}
\newtheorem{axiom}[theorem]{Axiom}
\theoremstyle{definition}
\newtheorem{definition}[theorem]{Definition}
\theoremstyle{remark}
\newtheorem{remark}[theorem]{Bemerkung}
\newtheorem{example}[theorem]{Beispiel}

\usepackage[unicode=true,hyperfootnotes=false]{hyperref}
\pdfstringdefDisableCommands{%
    \def\\{ }%
}
\hypersetup{
    pdftitle={Arithmetische Positivität und absolute Hodge-Klassen: Warum Hodge-Riemann-Positivität Delignes Absolutheit nicht verstärkt},
    pdfauthor={Lukas Geiger},
    pdfsubject={Grenzresultat zu Hodge-Riemann-Positivität und absoluten Hodge-Klassen},
    pdfkeywords={Hodge-Vermutung, absolute Hodge-Klassen, arithmetische Positivität, Hodge-Riemann-Relationen, Deligne-Absolutheit},
    colorlinks=true,
    linkcolor=black,
    urlcolor=blue,
    citecolor=black
}
\renewcommand*{\HyperDestNameFilter}[1]{ger.#1}
\renewcommand{\proofname}{Beweis}
\makeatletter
\let\auto@bib\@empty
\makeatother

\begin{document}

% ===================================================================
% TITELSEITE
% ===================================================================

\title{Arithmetische Positivität und absolute Hodge-Klassen:\\Warum Hodge-Riemann-Positivität\\Delignes Absolutheit nicht verstärkt}

\author{Lukas Geiger}
\affiliation{Unabhängiger Forscher, Bernau, Deutschland}

\date[Stand: ]{\today}

\begin{abstract}
Wir führen das Konzept der \textit{arithmetischen Positivität} für die Hodge-Vermutung ein: eine rationale $(p,p)$-Klasse $\alpha \in H^{2p}(X, \mathbb{Q})$ auf einer glatten projektiven Varietät $X/\mathbb{C}$ heißt arithmetisch positiv, wenn sie gleichzeitig Hodge-Riemann-Positivität unter allen Polarisierungen, Delignes Absolutheit unter allen Automorphismen von $\mathbb{C}$ und Lefschetz-Kompatibilität erfüllt. Wir beweisen, dass algebraische Klassen arithmetisch positiv sind, und dokumentieren begrenzte Funktorialitätstests: Der Rückzug entlang endlicher Morphismen ist auf dem durch Rückzug erzeugten Unterkegel des Amplekegels kontrolliert und wird nur unter einer expliziten Amplekegel-Vergleichshypothese zu einer vollen AP-Aussage; Künneth-Produkte werden nur als primitive Produktpolarisierungs-Vorzeichentests behandelt, sofern man nicht die spätere Identität $\mathrm{AP}=\mathrm{AbsHodge}$ nutzt; endlicher Vorwärtsschub bleibt eine Vermutung.

Die zentrale Beobachtung ist ein \textbf{Grenzresultat}: Der arithmetisch positive Kegel stimmt exakt mit Delignes absoluten Hodge-Klassen überein ($\mathrm{AP}^p(X) = \mathrm{AbsHodge}^p(X)$), weil die Hodge-Riemann-Bilinearrelationen die Positivitätsbedingungen für jede Klasse vom Typ $(p,p)$ automatisch garantieren. Dies ist eine wohlbekannte Konsequenz der Hodge-Riemann-Relationen; der Beitrag dieser Arbeit besteht darin, diese Sättigungsgrenze explizit zu machen und abzugrenzen, was Hodge-Riemann-basierte Positivitätsmethoden nicht leisten können. Der natürliche Versuch, Delignes Absolutheit durch Hodge-Riemann-Positivität zu verstärken, liefert daher keine neuen Einschränkungen. Die Arithmetische Positivitätsvermutung reduziert sich auf Delignes Frage (1982): ,,Sind alle absoluten Hodge-Klassen algebraisch?'' Ferner sind die Obstruktionen (O1) und (O3) in der No-Go-Formulierung für rationale $(p,p)$-Klassen vakuos; nur (O2), Versagen der Absolutheit, ist effektiv. Wir führen das Galois-Hodge-Riemann-Spektrum als Visualisierung des Absolutheitstests ein und diskutieren kontrollierte Testfamilien für Bedingungen, die genuin stärker als die Hodge-Riemann-Relationen sein müssten.

Diese Arbeit ist ein expositorischer Beitrag, der eine strukturelle Grenze dokumentiert; sie enthält kein neues Algebraizitätstheorem über die Beobachtung $\mathrm{AP} = \mathrm{AbsHodge}$ hinaus.
\end{abstract}

\keywords{Hodge-Vermutung, algebraische Zykel, arithmetische Positivität, absolute Hodge-Klassen, Hodge-Riemann-Bilinearform, No-Go-Satz}

\maketitle

\begin{center}
\footnotesize
ORCID: \url{https://orcid.org/0009-0005-7296-1534}\\
\end{center}

\section*{KI-Offenlegung}
KI-basierte Werkzeuge (Claude von Anthropic) wurden f\"ur mathematische Formalisierung, Literaturverifikation und redaktionelle \"Uberarbeitung eingesetzt. Der Autor konzipierte die Forschungsrichtung, pr\"ufte die mathematischen Inhalte und tr\"agt die volle Verantwortung f\"ur das Manuskript.

\clearpage

% ===================================================================
% 1. EINFUEHRUNG
% ===================================================================
\section{Einführung}
\label{sec:intro}

\subsection{Die Hodge-Vermutung}

Sei $X$ eine glatte projektive Varietät der Dimension $n$ über $\mathbb{C}$. Die Hodge-Zerlegung
\begin{equation}
H^k(X, \mathbb{C}) = \bigoplus_{p+q=k} H^{p,q}(X)
\label{eq:hodge_decomp}
\end{equation}
drückt die $k$-te singuläre Kohomologie als direkte Summe von Dolbeault-Kohomologiegruppen aus. Eine Klasse $\alpha \in H^{2p}(X, \mathbb{Q})$ heißt \textit{rationale Hodge-Klasse}, wenn ihr Bild in $H^{2p}(X, \mathbb{C})$ vollständig in $H^{p,p}(X)$ liegt. Jeder algebraische Zykel $Z$ der Kodimension $p$ definiert eine rationale Hodge-Klasse $\mathrm{cl}(Z) \in H^{2p}(X, \mathbb{Q}) \cap H^{p,p}(X)$ über die Zykelklassenabbildung.

\begin{conjecture}[Hodge, 1950 \cite{Hodge1950}]
\label{conj:hodge}
Jede rationale Hodge-Klasse auf einer glatten projektiven Varietät über $\mathbb{C}$ ist eine $\mathbb{Q}$-Linearkombination von Klassen algebraischer Zykel.
\end{conjecture}

Trotz über sieben Jahrzehnten der Forschung bleibt die Vermutung im Allgemeinen offen. Sie ist bekannt für:
\begin{itemize}
\item Divisoren ($p = 1$): durch den Lefschetz-$(1,1)$-Satz.
\item Abelsche Varietäten der Dimension $\leq 5$: durch Arbeiten von Hazama~\cite{Hazama1994}, Moonen~\cite{Moonen1999} und anderen.
\item Produkte elliptischer Kurven und bestimmte Fermat-Hyperflächen.
\item Uniruled-Varietäten (für Grad-$2$-Klassen).
\end{itemize}

\subsection{Grenzen rein konstruktiver Strategien}

Viele Ansätze zur Hodge-Vermutung versuchen, den algebraischen Zykel, der eine gegebene Hodge-Klasse repräsentiert, zu \textit{konstruieren}. Solche Konstruktionen sind in bekannten Spezialfällen unverzichtbar, zeigen aber zugleich strukturelle Engstellen. Diese Arbeit prüft eine komplementäre Route: ob nichtalgebraische Hodge-Klassen über eine Positivitätsobstruktion ausgeschlossen werden können.

Die für diese Route relevanten strukturellen Engstellen sind:

\textbf{1. Direkte Algebraisierung.} Gegeben eine rationale $(p,p)$-Klasse, versuche eine Untervarietät zu konstruieren, deren Fundamentalklasse mit ihr übereinstimmt. Dies scheitert, weil es keine kanonische Projektion von Hodge-Kohomologie auf algebraische Zykel gibt und die Konstruktion unter Deformation nicht funktoriell ist.

\textbf{2. Variationen von Hodge-Strukturen.} Verwende Familien $X_t$ und verfolge, wie sich Hodge-Klassen ändern. Dies scheitert, weil Hodge-Klassen ,,springen'' können -- die Bedingung, vom Typ $(p,p)$ zu sein, ist nicht offen in der Zariski-Topologie -- und Monodromie naive Stabilitätsargumente zerstört \cite{Voisin2002}.

\textbf{3. Motivische Methoden.} Ersetze Varietäten durch Motive und hoffe, dass Hodge-Klassen motivisch sind. Dies scheitert, weil die Kategorie der Motive nicht vollständig konstruiert ist: Grothendiecks Standardvermutungen \cite{Grothendieck1969} bleiben offen, und die Beziehung zwischen Hodge-Klassen und absoluten Hodge-Klassen ist ungeklärt \cite{Andre1996}.

\textbf{4. Reduktion mod $p$.} Reduziere auf Charakteristik $p$, berufe die Tate-Vermutung und lifte zurück. Dies scheitert, weil die Brücke Hodge $\Rightarrow$ Tate $\Rightarrow$ Hodge nicht geschlossen ist: Frobenius-Invarianz impliziert nicht den Hodge-Typ, und das Liften algebraischer Zykel ist nicht eindeutig \cite{Charles2013}.

Wie wir unten zeigen, stößt auch die in dieser Arbeit vorgeschlagene Positivitätsalternative auf eine fundamentale Einschränkung: die natürlichen Positivitätsbedingungen reduzieren sich auf Delignes Absolutheit. Dennoch markiert die Analyse eine präzise Grenze für diese Klasse von Ansätzen.

\subsection{Die Positivitätsalternative}

Wir beobachten, dass große Durchbrüche in benachbarten Vermutungen einem anderen Muster folgen:

\begin{center}
\resizebox{\columnwidth}{!}{%
\begin{tabular}{ll}
\toprule
Vermutung & Zentrale Positivitätseigenschaft \\
\midrule
Weil-Vermutungen & Positivität der Frobenius-Eigenwerte \\
BSD (Rang 1) & Positivität der Höhen (Gross--Zagier) \\
Tate (abelsche Var.) & Positivität via Faltings' Satz \\
\bottomrule
\end{tabular}}
\end{center}

In jedem Fall spielt eine Positivitätseigenschaft des jeweiligen arithmetischen Objekts eine strukturelle Rolle, auch wenn sich die konkreten Beweistechniken deutlich unterscheiden. Der Hodge-Vermutung fehlt genau dies: \textit{eine Positivitätsstruktur, die Algebraizität erzwingen könnte, ohne Zykel explizit zu konstruieren}.

Diese Arbeit führt eine solche Struktur ein -- \textit{arithmetische Positivität} -- die Hodge-Riemann-Positivität unter allen Polarisierungen mit Delignes Absolutheit unter allen Automorphismen von $\mathbb{C}$ kombiniert. Die zentrale Beobachtung ist jedoch ein Grenz-/No-Go-Resultat: die Hodge-Riemann-Bilinearrelationen garantieren die Positivitätsbedingungen für jede Klasse vom Typ $(p,p)$ automatisch, sodass der arithmetisch positive Kegel mit Delignes absoluten Hodge-Klassen übereinstimmt: $\mathrm{AP}^p(X) = \mathrm{AbsHodge}^p(X)$. Diese automatische Sättigung ist Spezialisten bekannt; Ziel dieser Arbeit ist es, die Sättigungsgrenze in einem einheitlichen Positivitätsrahmen explizit zu machen. Der Versuch, Absolutheit durch Hodge-Riemann-Positivität zu verstärken, liefert keine neuen Einschränkungen. Die Arithmetische Positivitätsvermutung reduziert sich auf Delignes Frage: ,,Sind alle absoluten Hodge-Klassen algebraisch?''

Wir betrachten dieses negative Resultat als informativ: es markiert eine präzise Grenze für positivitätsbasierte Ansätze zur Hodge-Vermutung und identifiziert das Galois-Hodge-Riemann-Spektrum als diagnostisches Werkzeug.

\subsection{Gliederung}

Abschnitt~\ref{sec:background} wiederholt die Hodge-Riemann-Bilinearform und absolute Hodge-Klassen. Abschnitt~\ref{sec:arithmetic_positivity} führt arithmetische Positivität ein. Abschnitt~\ref{sec:easy_direction} beweist, dass algebraische Klassen arithmetisch positiv sind. Abschnitt~\ref{sec:functoriality} dokumentiert begrenzte Funktorialitätstests und ihre Hypothesen. Abschnitt~\ref{sec:nogo} formuliert den No-Go-Satz. Abschnitt~\ref{sec:abelian} beweist die Äquivalenz für abelsche Varietäten. Abschnitt~\ref{sec:obstruction} identifiziert die präzise Obstruktion. Abschnitt~\ref{sec:discussion} diskutiert die verbleibende Lücke und die Beziehung zu Delignes Programm.


% ===================================================================
% 2. HINTERGRUND
% ===================================================================
\section{Hintergrund}
\label{sec:background}

\subsection{Hodge-Riemann-Bilinearform}

Sei $X$ eine glatte projektive Varietät der Dimension $n$, ausgestattet mit einer K\"ahler-Form $\omega$, die eine ample Klasse $L = [\omega] \in H^{1,1}(X) \cap H^2(X, \mathbb{Z})$ definiert. Der \textit{Lefschetz-Operator} $\Lambda_L: H^k(X) \to H^{k+2}(X)$ ist definiert durch $\Lambda_L(\alpha) = L \cup \alpha$. Wir verwenden die Konvention $\Lambda_L(\alpha) = L \cup \alpha$; einige Autoren reservieren $\Lambda$ für den dualen Lefschetz-Operator und notieren die Cup-Produkt-Operation einfach als $L$ oder $L \wedge$ (Voisin~\cite{Voisin2002}; Griffiths--Harris~\cite{GH1978}).

\begin{theorem}[Hard Lefschetz]
\label{thm:hard_lefschetz}
Für $k \leq n$ ist die Abbildung $\Lambda_L^{n-k}: H^k(X, \mathbb{Q}) \to H^{2n-k}(X, \mathbb{Q})$ ein Isomorphismus.
\end{theorem}

Eine Klasse $\alpha \in H^k(X, \mathbb{Q})$ mit $k \leq n$ ist \textit{primitiv} (bezüglich $L$), wenn $\Lambda_L^{n-k+1}(\alpha) = 0$. Die \textit{Hodge-Riemann-Bilinearform} auf der primitiven Kohomologie $P^k(X, \mathbb{Q}) \subset H^k(X, \mathbb{Q})$ ist:
\begin{equation}
Q_L(\alpha, \beta) = \int_X \alpha \cup \beta \cup L^{n-k}\,.
\label{eq:HR_form}
\end{equation}

\begin{theorem}[Hodge-Riemann-Bilinearrelationen]
\label{thm:HR}
Auf dem primitiven Unterraum $P^{p,q}(X) = P^{p+q}(X) \cap H^{p,q}(X)$ mit $p + q = k$ ist die hermitesche Form
\begin{equation}
h_L(\alpha, \beta) = i^{p-q}\,(-1)^{k(k-1)/2}\, Q_L(\alpha, \bar{\beta})
\label{eq:HR_hermitian}
\end{equation}
positiv definit.
\end{theorem}

Die Hodge-Riemann-Relationen sind das fundamentale Positivitätsresultat in der Hodge-Theorie. Sie liefern Definitheit \textit{relativ zu einer festen Polarisierung} $L$. Ein zentrales Thema dieser Arbeit ist, dass Algebraizität Positivität nicht nur für ein $L$ erfordert, sondern für \textit{alle} zulässigen Polarisierungen und Operationen.

\subsection{Absolute Hodge-Klassen}

Deligne \cite{Deligne1982} führte eine Verfeinerung von Hodge-Klassen ein. Ein Automorphismus $\sigma \in \mathrm{Aut}(\mathbb{C})$ sendet eine glatte projektive Varietät $X/\mathbb{C}$ auf die konjugierte Varietät $X^\sigma$ (gleiche Gleichungen, Koeffizienten verdrillt durch $\sigma$). Der Vergleichsisomorphismus liefert:
\begin{equation}
H^{2p}_{\mathrm{dR}}(X/\mathbb{C}) \xrightarrow{\sim} H^{2p}_{\mathrm{dR}}(X^\sigma/\mathbb{C})\,,
\end{equation}
der die de~Rham-Klasse von $\alpha$ auf $X$ auf eine Klasse $\alpha^\sigma$ auf $X^\sigma$ abbildet.

\begin{definition}[Deligne]
\label{def:absolute}
Eine rationale Hodge-Klasse $\alpha \in H^{2p}(X, \mathbb{Q}) \cap H^{p,p}(X)$ ist \textit{absolut} (oder \textit{absolut Hodge}), wenn für jeden $\sigma \in \mathrm{Aut}(\mathbb{C})$ die konjugierte Klasse $\alpha^\sigma$ wieder eine Hodge-Klasse auf $X^\sigma$ ist.
\end{definition}

\begin{theorem}[Deligne \cite{Deligne1982}]
\label{thm:deligne_absolute}
\begin{enumerate}
\item Jede algebraische Klasse ist absolut.
\item Auf abelschen Varietäten ist jede Hodge-Klasse absolut.
\end{enumerate}
\end{theorem}

Delignes Resultat zeigt, dass Absolutheit eine \textit{notwendige} Bedingung für Algebraizität ist. Die Frage ist, ob sie hinreichend ist. Unser Konzept der arithmetischen Positivität \textit{versucht}, Absolutheit durch Hinzufügen von Hodge-Riemann-Positivität unter allen Konjugationen zu verstärken; wie in Abschnitt~\ref{sec:discussion} gezeigt, gelingt dies nicht (die zusätzlichen Bedingungen sind automatisch), aber die Analyse ist informativ.


% ===================================================================
% 3. ARITHMETISCHE POSITIVITAET
% ===================================================================
\section{Arithmetische Positivität}
\label{sec:arithmetic_positivity}

Wir führen nun das zentrale Konzept ein.

\subsection{Definition}

Sei $X$ eine glatte projektive Varietät der Dimension $n$ über $\mathbb{C}$, und sei $\mathrm{Amp}(X)$ der Amplekegel von $X$.

\begin{definition}[Arithmetische Positivität]
\label{def:arith_pos}
Eine rationale Hodge-Klasse $\alpha \in H^{2p}(X, \mathbb{Q}) \cap H^{p,p}(X)$ ist \textit{arithmetisch positiv}, wenn sie die folgenden drei Bedingungen gleichzeitig erfüllt:

In der direkten primitiven Notation unten arbeiten wir im mittleren oder unteren Bereich $2p\leq n$; Klassen in komplementärer Kodimension sind über Poincar\'e-Dualität zu verstehen.  Entsprechend wird die Lefschetz-Kompatibilität nur für Cup-Potenzen gefordert, deren primitive Komponente im mittleren oder unteren Bereich direkt definiert ist.

\begin{enumerate}
\item[\textbf{(AP1)}] \textbf{Universelle Hodge-Riemann-Positivität.} Für jede ample Klasse $L \in \mathrm{Amp}(X)$ erfüllt die primitive Komponente $\alpha_0$ von $\alpha$ (bezüglich $L$) die \textit{vorzeichenbehaftete} Hodge-Riemann-Ungleichung:
\begin{equation}
(-1)^p\, Q_L(\alpha_0, \alpha_0) \geq 0\,,
\label{eq:AP1}
\end{equation}
wobei $Q_L$ die Hodge-Riemann-Form~\eqref{eq:HR_form} ist. Das Vorzeichen $(-1)^p$ ist wesentlich: die Hodge-Riemann-Bilinearrelationen (Satz~\ref{thm:HR}) ergeben $(-1)^p Q_L > 0$ auf primitiven $(p,p)$-Klassen, nicht $Q_L > 0$.

\item[\textbf{(AP2)}] \textbf{Absolute Stabilität.} Für jeden $\sigma \in \mathrm{Aut}(\mathbb{C})$:
\begin{enumerate}
\item[(a)] $\alpha^\sigma$ ist eine Hodge-Klasse auf $X^\sigma$ (d.h., $\alpha$ ist absolut).
\item[(b)] Die konjugierte Klasse $\alpha^\sigma$ erfüllt (AP1) auf $X^\sigma$ bezüglich jeder amplen Klasse $L^\sigma \in \mathrm{Amp}(X^\sigma)$.
\end{enumerate}

\item[\textbf{(AP3)}] \textbf{Lefschetz-Kompatibilität.} Für alle $0 \leq j \leq \lfloor n/2\rfloor - p$ und alle amplen $L$, sei $\gamma_j = \Lambda_L^j \alpha \in H^{2(p+j)}(X)$ und $\gamma_{j,0}$ ihre primitive Komponente bezüglich $L$ in $H^{2(p+j)}(X)$. Dann:
\begin{equation}
(-1)^{p+j}\, Q_{L}(\gamma_{j,0},\, \gamma_{j,0}) \geq 0\,.
\label{eq:AP3}
\end{equation}
\end{enumerate}

Die Menge der arithmetisch positiven Klassen in $H^{2p}(X, \mathbb{Q}) \cap H^{p,p}(X)$ wird mit $\mathrm{AP}^p(X)$ bezeichnet.
\end{definition}

\begin{remark}
\label{rem:redundancy}
Die Bedingungen (AP1), (AP2b) und (AP3) sind \textit{nicht} unabhängig von (AP2a). Wie in Proposition~\ref{prop:absolute_implies_ap} unten gezeigt, garantieren die Hodge-Riemann-Bilinearrelationen (AP1), (AP2b) und (AP3) automatisch für jede Klasse vom Typ $(p,p)$. Die einzige substantielle Bedingung ist daher (AP2a): Delignes Absolutheit. Insbesondere gilt $\mathrm{AP}^p(X) = \mathrm{AbsHodge}^p(X)$ (Bemerkung~\ref{rem:ap_equals_absolute}). Die Definition behält alle drei Bedingungen aus darstellerischen Gründen bei: sie machen den \textit{Typ} der Positivität explizit, den algebraische Klassen genießen, auch wenn die Hodge-Riemann-Relationen sie für $(p,p)$-Klassen automatisch machen. Genuin stärkere Bedingungen zu finden -- jenseits der Hodge-Riemann-Relationen -- die \textit{tatsächlich} über Absolutheit hinaus einschränken, bleibt ein offenes Problem (siehe Abschnitt~\ref{sec:discussion}).
\end{remark}

\subsection{Die Arithmetische Positivitätsvermutung}

\begin{conjecture}[Arithmetische Positivitätsvermutung]
\label{conj:AP}
Sei $X$ eine glatte projektive Varietät über $\mathbb{C}$. Dann gilt:
\begin{equation}
\mathrm{AP}^p(X) = \mathrm{cl}\!\left(\mathrm{CH}^p(X)_\mathbb{Q}\right)\,,
\end{equation}
wobei $\mathrm{CH}^p(X)_\mathbb{Q}$ die Chow-Gruppe der Kodimension-$p$-Zykel mit rationalen Koeffizienten ist und $\mathrm{cl}$ die Zykelklassenabbildung.
\end{conjecture}

\begin{proposition}
\label{prop:equivalence}
Die folgenden Beziehungen gelten zwischen der APV und der HV:
\begin{enumerate}
\item Die Hodge-Vermutung impliziert die Arithmetische Positivitätsvermutung.
\item Die ,,schwere Richtung'' der APV ($\mathrm{AP}^p(X) \subseteq \mathrm{cl}(\mathrm{CH}^p(X)_\mathbb{Q})$) zusammen mit der Aussage, dass alle rationalen Hodge-Klassen arithmetisch positiv sind ($H^{2p}(X, \mathbb{Q}) \cap H^{p,p}(X) \subseteq \mathrm{AP}^p(X)$), impliziert die HV.
\item Unbedingt: $\mathrm{cl}(\mathrm{CH}^p(X)_\mathbb{Q}) \subseteq \mathrm{AP}^p(X)$ (Satz~\ref{thm:easy_direction} unten).
\end{enumerate}
\end{proposition}

\begin{proof}
(1) Wenn die HV gilt, ist jede Hodge-Klasse algebraisch, und algebraische Klassen erfüllen (AP1)--(AP3) (Satz~\ref{thm:easy_direction} unten). Somit $\mathrm{AP}^p(X) \subseteq H^{2p}(X, \mathbb{Q}) \cap H^{p,p}(X) = \mathrm{cl}(\mathrm{CH}^p(X)_\mathbb{Q})$, und die umgekehrte Inklusion gilt nach Satz~\ref{thm:easy_direction}, also $\mathrm{AP}^p(X) = \mathrm{cl}(\mathrm{CH}^p(X)_\mathbb{Q})$.

(2) Dies erfordert zwei separate Annahmen. \textit{Erstens}, die schwere Richtung der APV gelte: $\mathrm{AP}^p(X) \subseteq \mathrm{cl}(\mathrm{CH}^p(X)_\mathbb{Q})$, d.h.\ jede arithmetisch positive Klasse ist algebraisch. \textit{Zweitens}, jede rationale Hodge-Klasse sei arithmetisch positiv: $H^{2p}(X, \mathbb{Q}) \cap H^{p,p}(X) \subseteq \mathrm{AP}^p(X)$. Unter beiden Annahmen liegt jede rationale Hodge-Klasse $\alpha$ in $\mathrm{AP}^p(X)$ (nach der zweiten Annahme) und damit in $\mathrm{cl}(\mathrm{CH}^p(X)_\mathbb{Q})$ (nach der ersten Annahme), also ist $\alpha$ algebraisch. Dies ergibt die HV.

Keine der beiden Annahmen genügt allein: die APV ($\mathrm{AP}^p = \mathrm{alg}^p$) zeigt, dass AP-Klassen algebraisch sind, aber ohne das Wissen, dass alle Hodge-Klassen AP sind, kann man nicht schließen, dass alle Hodge-Klassen algebraisch sind.

(3) Siehe Satz~\ref{thm:easy_direction}.
\end{proof}

\begin{remark}
Die APV allein impliziert die HV nicht direkt. Die APV behauptet $\mathrm{AP}^p = \mathrm{alg}^p$, während die HV $\mathrm{Hodge}^p = \mathrm{alg}^p$ behauptet. Da $\mathrm{AP}^p \subseteq \mathrm{Hodge}^p$ (arithmetische Positivität stellt zusätzliche Bedingungen jenseits einer Hodge-Klasse), gibt die APV $\mathrm{alg}^p \subseteq \mathrm{Hodge}^p$ (trivial), aber nicht $\mathrm{Hodge}^p \subseteq \mathrm{alg}^p$ ohne das zusätzliche Wissen, dass alle Hodge-Klassen arithmetisch positiv sind. Der nichttriviale Inhalt zerfällt in zwei Teile: (i)~zeige dass nicht-AP Hodge-Klassen nicht existieren (d.h. jede Hodge-Klasse ist AP), und (ii)~zeige dass AP Algebraizität impliziert (die schwere Richtung der APV).
\end{remark}

Der nichttriviale Inhalt der APV ist die ,,schwere Richtung'': $\mathrm{AP}^p(X) \subseteq \mathrm{cl}(\mathrm{CH}^p(X)_\mathbb{Q})$. Diese besagt, dass arithmetische Positivität Algebraizität \textit{erzwingt}.


% ===================================================================
% 4. DIE LEICHTE RICHTUNG
% ===================================================================
\section{Algebraische Klassen sind arithmetisch positiv}
\label{sec:easy_direction}

\begin{theorem}[Leichte Richtung]
\label{thm:easy_direction}
Sei $X$ eine glatte projektive Varietät über $\mathbb{C}$. Jede Klasse im Bild der Zykelklassenabbildung ist arithmetisch positiv:
\begin{equation}
\mathrm{cl}\!\left(\mathrm{CH}^p(X)_\mathbb{Q}\right) \subseteq \mathrm{AP}^p(X)\,.
\end{equation}
\end{theorem}

\begin{proof}
Sei $\alpha = \mathrm{cl}(Z)$ für einen algebraischen Zykel $Z$ der Kodimension $p$ mit rationalen Koeffizienten. Wir verifizieren (AP1)--(AP3).

\textit{(AP1):} Sei $\alpha_0$ die $L$-primitive Komponente von $\alpha = \mathrm{cl}(Z)$. Da $\alpha$ algebraisch ist, liegt es in $H^{p,p}(X) \cap H^{2p}(X, \mathbb{Q})$, und seine primitive Komponente $\alpha_0$ liegt in $P^{p,p}(X) \cap H^{2p}(X, \mathbb{Q})$. Nach den Hodge-Riemann-Bilinearrelationen (Satz~\ref{thm:HR}) ist die hermitesche Form $h_L$ positiv definit auf $P^{p,p}(X)$. Für eine rationale $(p,p)$-Klasse gilt $\bar{\alpha}_0 = \alpha_0$, $i^{p-p} = 1$ und $(-1)^{k(k-1)/2} = (-1)^{p(2p-1)}$ mit $k = 2p$. Also $h_L(\alpha_0, \alpha_0) = (-1)^{p(2p-1)}\, Q_L(\alpha_0, \alpha_0)$. Da $(-1)^{p(2p-1)} = (-1)^{2p^2 - p} = (-1)^p$ (weil $(-1)^{2p^2} = 1$), ergeben die Hodge-Riemann-Relationen:
\begin{equation}
(-1)^p\, Q_L(\alpha_0, \alpha_0) > 0
\end{equation}
falls $\alpha_0 \neq 0$. Für $\alpha_0 = 0$ (d.h., wenn $\alpha$ keine primitive Komponente hat) gilt die Ungleichung trivial. Dieses Argument verwendet nur die Hodge-Riemann-Relationen auf $P^{p,p}(X)$ und erfordert nicht, dass $\alpha_0$ selbst die Klasse eines effektiven Zykels ist.

\textit{(AP2):} Teil (a) ist Delignes Satz (Satz~\ref{thm:deligne_absolute}): algebraische Klassen sind absolut. Teil (b) folgt, weil $\sigma$ algebraische Zykel auf $X$ auf algebraische Zykel auf $X^\sigma$ sendet (Algebraizität ist eine Eigenschaft polynomialer Gleichungen, die kovariant unter $\sigma$ transformieren), und algebraische Klassen auf $X^\sigma$ (AP1) durch das obige Argument erfüllen.

\textit{(AP3):} Cup-Produkt mit einer amplen $(1,1)$-Klasse erhält den Hodge-Typ: Für zulässige $j$ in Definition~\ref{def:arith_pos} liegt $\gamma_j=\Lambda_L^j\alpha=L^j\cup\alpha$ in $H^{p+j,p+j}(X)$. Ist $L$ integral oder rational, kann dies als Schnitt von $Z$ mit Hyperebenenklassen repräsentiert werden; für eine allgemeine reelle ample Klasse ist die Zykelsprache nur nach rationaler Approximation zu lesen, während das Vorzeichenargument selbst nur den Hodge-Typ benutzt. Die $L$-primitive Komponente $\gamma_{j,0}$ liegt in $P^{p+j,\,p+j}(X)$, und die Hodge-Riemann-Relationen liefern $(-1)^{p+j} Q_L(\gamma_{j,0}, \gamma_{j,0}) \geq 0$, mit demselben Argument wie in (AP1).
\end{proof}


% ===================================================================
% 5. FUNKTORIALITAET
% ===================================================================
\section{Funktorialität der Arithmetischen Positivität}
\label{sec:functoriality}

Für ein vollständiges No-Go-Programm müsste der arithmetisch positive Kegel unter den Standardoperationen der algebraischen Geometrie funktoriell sein. Die folgenden Resultate dokumentieren nur die begrenzte Funktorialität, die hier tatsächlich kontrolliert wird.

\begin{proposition}[Rückzug -- endliche Morphismen]
\label{prop:pullback}
Sei $f: Y \to X$ ein endlicher Morphismus glatter projektiver Varietäten. Wenn $\alpha \in \mathrm{AP}^p(X)$, dann erfüllt $f^*\alpha$ (AP1) für ample Klassen im Rückzugsunterkegel $f^*\mathrm{Amp}(X) \subset \mathrm{Amp}(Y)$ und erfüllt (AP2)--(AP3) auf diesem Unterkegel. Ein volles Element von $\mathrm{AP}^p(Y)$ erhält man nur unter einer expliziten Amplekegel-Vergleichshypothese, etwa wenn die relevanten Positivitätsungleichungen durch Dichtheit oder Gleichheit der numerischen Amplekegel von $f^*\mathrm{Amp}(X)$ auf ganz $\mathrm{Amp}(Y)$ fortgesetzt werden können.
\end{proposition}

\begin{proof}
Sei $f: Y \to X$ endlich vom Grad $d$.

\textit{(AP1):} Da $f$ endlich ist, ist $f^*L_X$ ampel auf $Y$ für jede ample Klasse $L_X \in \mathrm{Amp}(X)$. Für die primitive Komponente $\alpha_0$ von $\alpha$ bezüglich $L_X$ liefert die Projektionsformel:
\[
\int_Y f^*\alpha_0 \cup f^*\bar{\alpha}_0 \cup (f^*L_X)^{n-2p} = d \int_X \alpha_0 \cup \bar{\alpha}_0 \cup L_X^{n-2p},
\]
wobei $n = \dim X = \dim Y$. Da $d > 0$ und die rechte Seite $(-1)^p Q_{L_X}(\alpha_0, \alpha_0) \geq 0$ nach Voraussetzung erfüllt, erhalten wir $(-1)^p Q_{f^*L_X}(f^*\alpha_0, f^*\alpha_0) \geq 0$ auf dem Rückzugsunterkegel. Für eine allgemeine ample Klasse $L_Y \in \mathrm{Amp}(Y)$ folgt aus Endlichkeit allein keine Fortsetzung: Die primitive Komponente von $f^*\alpha$ hängt von $L_Y$ ab, und der Rückzugsunterkegel muss nicht dicht im Amplekegel von $Y$ liegen. Wenn eine zusätzliche Amplekegel-Vergleichshypothese Dichtheit oder Gleichheit des relevanten Kegels liefert, setzen Stetigkeit der Hodge-Riemann-Form und der Lefschetz-Zerlegung die Ungleichung auf alle amplen Klassen fort.

\textit{(AP2):} Rückzug kommutiert mit $\sigma$-Konjugation: $(f^*\alpha)^\sigma = (f^\sigma)^*(\alpha^\sigma)$, und $f^\sigma: Y^\sigma \to X^\sigma$ ist wieder endlich vom selben Grad. Da $\alpha^\sigma \in \mathrm{AP}^p(X^\sigma)$ nach Voraussetzung, folgt das Resultat aus dem (AP1)-Argument angewandt auf $f^\sigma$ und $\alpha^\sigma$.

\textit{(AP3):} Der Lefschetz-Operator erfüllt $\Lambda_{f^*L}^j \circ f^* = f^* \circ \Lambda_L^j$ (Kompatibilität von Rückzug mit Cup-Produkt). Daher gilt $\Lambda_{f^*L}^j(f^*\alpha) = f^*(\Lambda_L^j \alpha)$, und die Positivität der primitiven Komponente folgt aus dem Grad-$d$-Argument angewandt auf $\Lambda_L^j\alpha$.
\end{proof}

\begin{remark}
Die im Beweis verwendete Gradformel $Q_{f^*L_X}(f^*\alpha_0, f^*\alpha_0) = \deg(f) \cdot Q_{L_X}(\alpha_0, \alpha_0)$ setzt voraus, dass $f$ ein endlicher Morphismus ist (äquivalent: generisch endlich mit $\dim Y = \dim X$). Für einen generisch endlichen Morphismus $f: Y \to X$ mit $\dim Y = \dim X$, der nicht endlich ist, liefert die Projektionsformel die Gradrelation zwar auf einer dichten offenen Teilmenge, aber der Verzweigungsort erfordert zusätzliche Sorgfalt; die Aussage von Proposition~\ref{prop:pullback} ist daher auf endliche Morphismen beschränkt.
\end{remark}

\begin{remark}
Für Morphismen $f: Y \to X$ mit $\dim Y > \dim X$ ist der Rückzug $f^*L_X$ nicht ampel auf $Y$, und die primitive Zerlegung auf $Y$ erfordert eine relative Hodge-Riemann-Analyse. Die Funktorialität von $\mathrm{AP}^p$ unter allgemeinen Morphismen bleibt offen.
\end{remark}

\begin{remark}[Dichte-Vorbehalt]\label{rem:density-caveat}
Die Erweiterung von $f^*\mathrm{Amp}(X)$ auf ganz $\mathrm{Amp}(Y)$ erfordert einen expliziten Vergleich der Amplekegel. Eine nützliche hinreichende Bedingung ist, dass der Rückzugsunterkegel ein dichtes Bild im relevanten numerischen Amplekegel hat oder dass die numerischen Amplekegel nach Rückzug übereinstimmen. Das kann in Picard-Zahl-eins-Situationen gelten, scheitert aber im Allgemeinen. Ohne solche Zusatzannahme ist Proposition~\ref{prop:pullback} nur eine Funktorialitätsaussage auf dem Rückzugsunterkegel.
\end{remark}

\begin{conjecture}[Vorwärtsschub -- endliche Morphismen]
\label{prop:pushforward}
Sei $f: Y \to X$ ein endlicher Morphismus glatter projektiver Varietäten (also $\dim Y = \dim X$). Wenn $\alpha \in \mathrm{AP}^{p}(Y)$, dann gilt $f_*\alpha \in \mathrm{AP}^p(X)$.
\end{conjecture}

\begin{remark}[Evidenzstatus für endlichen Vorwärtsschub]
Die folgende Heuristik stützt Vermutung~\ref{prop:pushforward}; sie ist kein
Beweis. Die Projektionsformel $f_*f^* = \deg(f) \cdot \mathrm{id}$ auf
$H^*(X, \mathbb{Q})$ impliziert, dass $f_*$ auf dem Bild von $f^*$ injektiv ist
(bis auf Skalar). Für $\alpha \in \mathrm{AP}^p(Y)$ und jede ample Klasse
$L_X \in \mathrm{Amp}(X)$ ist $f^*L_X$ ampel auf $Y$ (da $f$ endlich). Die
Hodge-Riemann-Form auf $X$ ist mit jener auf $Y$ durch die Adjunktion
$\int_X f_*\beta \cup \gamma = \int_Y \beta \cup f^*\gamma$ verknüpft.
Zusammen mit der Gradformel liefert dies Evidenz für die Vorzeichenbedingungen
von $f_*\alpha$. Die absolute Stabilität sollte sich übertragen, da $f^\sigma$
wieder endlich vom selben Grad ist. Ein vollständiger Beweis erfordert die
Verifikation, dass die primitive Zerlegung auf $X$ kompatibel mit $f_*$
angewandt auf primitive Klassen auf $Y$ ist; dies ist nicht offensichtlich, da
$f_*$ Primitivität im Allgemeinen nicht erhält.
\end{remark}

\begin{remark}
Für Morphismen $f: Y \to X$ mit $\dim Y > \dim X$ erfordert der Vorwärtsschub $f_*$ ein Gysin-Argument, das die Hodge-Riemann-Form auf $Y$ mit jener auf $X$ über die Grothendieck-Riemann-Roch-Formel und die relative primitive Zerlegung in Beziehung setzt. Diese Analyse wird hier nicht durchgeführt; die Funktorialität von $\mathrm{AP}^p$ unter allgemeinem Vorwärtsschub bleibt offen.
\end{remark}

\begin{proposition}[Künneth-Produkt -- primitiver Produktpolarisierungs-Vorzeichentest]
\label{prop:kunneth}
Seien $\alpha \in \mathrm{AP}^p(X)$ und $\beta \in \mathrm{AP}^q(Y)$ primitiv bezüglich amplen Klassen $L_X$ bzw.\ $L_Y$. Dann beweist die direkte Produktform-Rechnung die (AP1)-Vorzeichenungleichung für $\alpha \boxtimes \beta$ für alle Produktpolarisierungen $L = L_X \boxtimes 1 + 1 \boxtimes L_Y$, und Konjugation ist mit dem äußeren Produkt kompatibel. Diese Proposition ist nur ein Produktpolarisierungs-Vorzeichentest; sie beweist nicht aus sich heraus die volle (AP3)-Bedingung für alle Lefschetz-Potenzen auf $X\times Y$.
\end{proposition}

\begin{remark}
Die Obstruktion besteht nicht nur darin, dass nichtprimitive Eingangsklassen schwieriger sind. Auch wenn $\alpha$ und $\beta$ primitiv sind, zerlegen sich die Klassen $\Lambda_L^j(\alpha\boxtimes\beta)$ im Allgemeinen in Summen, deren $L$-primitive Projektionen auf $X\times Y$ Kreuzterme enthalten. Die Vorzeichenformel unten kontrolliert das primitive äußere Produkt selbst unter Produktpolarisierungen. Volle AP-Mitgliedschaft von $\alpha\boxtimes\beta$ wird in diesem Paper nur über die spätere Sättigungsidentität $\mathrm{AP}=\mathrm{AbsHodge}$ erhalten, zusammen mit der Standardtatsache, dass äußere Produkte absoluter Hodge-Klassen absolut sind; sie ist kein eigenständiger Künneth-AP-Satz.
\end{remark}

Wir etablieren zunächst ein Schlüssellemma über die Primitivität von äußeren Produkten.

\begin{lemma}[Primitivität äußerer Produkte]
\label{lem:primitive_product}
Seien $X$ und $Y$ glatte projektive Varietäten der Dimensionen $n$ bzw.\ $m$, mit amplen Klassen $L_X \in \mathrm{Amp}(X)$ und $L_Y \in \mathrm{Amp}(Y)$. Setze $L = L_X \boxtimes 1 + 1 \boxtimes L_Y \in \mathrm{Amp}(X \times Y)$. Wenn $\alpha \in P^{2p}(X, \mathbb{Q})$ primitiv bezüglich $L_X$ und $\beta \in P^{2q}(Y, \mathbb{Q})$ primitiv bezüglich $L_Y$ ist, dann ist $\alpha \boxtimes \beta$ primitiv bezüglich $L$ auf $X \times Y$.
\end{lemma}

\begin{proof}
Setze $k = p + q$ und $N = n + m = \dim(X \times Y)$. Wir müssen zeigen, dass $L^{N-2k+1} \cup (\alpha \boxtimes \beta) = 0$. Da $L = L_X \boxtimes 1 + 1 \boxtimes L_Y$, liefert die Binomialentwicklung für beliebiges $j \geq 0$:
\[
L^j \cup (\alpha \boxtimes \beta) = \sum_{a+b=j} \binom{j}{a} (L_X^a \cup \alpha) \boxtimes (L_Y^b \cup \beta).
\]
Jeder Summand verschwindet, es sei denn beide Faktoren sind von null verschieden. Da $\alpha$ $L_X$-primitiv ist, gilt $L_X^{n-2p+1} \cup \alpha = 0$, also $L_X^a \cup \alpha = 0$ für $a \geq n - 2p + 1$. Analog gilt $L_Y^b \cup \beta = 0$ für $b \geq m - 2q + 1$. Ein Summand mit $a + b = j$ ist nur dann von null verschieden, wenn $a \leq n - 2p$ und $b \leq m - 2q$, was $j \leq (n - 2p) + (m - 2q) = N - 2k$ erfordert. Daher gilt $L^j \cup (\alpha \boxtimes \beta) = 0$ für alle $j \geq N - 2k + 1$, was genau die Primitivitätsbedingung ist.
\end{proof}

\begin{lemma}[Vorzeichenkompatibilität für äußere Produkte]
\label{lem:sign_product}
Unter den Voraussetzungen von Lemma~\ref{lem:primitive_product}, mit $k = p + q$:
\begin{equation}
\begin{aligned}
(-1)^k\, Q_L(\alpha \boxtimes \beta,\, \alpha \boxtimes \beta)
&= \binom{N-2k}{n-2p}
   \bigl[(-1)^p Q_{L_X}(\alpha, \alpha)\bigr]\\
&\quad\cdot \bigl[(-1)^q Q_{L_Y}(\beta, \beta)\bigr],
\end{aligned}
\label{eq:sign_kunneth}
\end{equation}
wobei $N = n + m$ und $Q_L$, $Q_{L_X}$, $Q_{L_Y}$ die Hodge-Riemann-Formen~\eqref{eq:HR_form} auf $X \times Y$, $X$ bzw.\ $Y$ sind.
\end{lemma}

\begin{proof}
Nach Lemma~\ref{lem:primitive_product} ist $\alpha \boxtimes \beta$ $L$-primitiv, also liefert die Definition~\eqref{eq:HR_form} direkt:
\[
Q_L(\alpha \boxtimes \beta,\, \alpha \boxtimes \beta) = \int_{X \times Y} (\alpha \boxtimes \beta)^2 \cup L^{N-2k}.
\]
Entwicklung von $L^{N-2k}$ über die Binomialformel wie in Lemma~\ref{lem:primitive_product} zeigt, dass der einzige nichtverschwindende Term bei $a = n - 2p$, $b = m - 2q$ auftritt (erzwungen durch $a \leq n-2p$, $b \leq m-2q$ und $a + b = N - 2k = (n-2p) + (m-2q)$). Nach der Künneth-Formel für Integration:
\[
\begin{aligned}
\int_{X \times Y}(\alpha \boxtimes \beta)^2 \cup L^{N-2k}
&= \binom{N-2k}{n-2p}
   \left(\int_X \alpha^2 \cup L_X^{n-2p}\right)\\
&\quad\cdot
   \left(\int_Y \beta^2 \cup L_Y^{m-2q}\right).
\end{aligned}
\]
Nach Definition~\eqref{eq:HR_form} sind die Integrale auf $X$ und $Y$ gerade $Q_{L_X}(\alpha, \alpha)$ bzw.\ $Q_{L_Y}(\beta, \beta)$. Daher:
\[
\begin{aligned}
Q_L(\alpha \boxtimes \beta,\, \alpha \boxtimes \beta)
&= \binom{N-2k}{n-2p}\,
   Q_{L_X}(\alpha, \alpha)\\
&\quad\cdot Q_{L_Y}(\beta, \beta).
\end{aligned}
\]
Multiplikation beider Seiten mit $(-1)^k = (-1)^{p+q} = (-1)^p \cdot (-1)^q$ ergibt Gleichung~\eqref{eq:sign_kunneth}.
\end{proof}

\begin{proof}[Beweis von Proposition~\ref{prop:kunneth}]
Wir verifizieren die kontrollierte (AP1)-Aussage für Produktpolarisierungen für $\alpha \boxtimes \beta$ mit $\alpha \in \mathrm{AP}^p(X)$ primitiv und $\beta \in \mathrm{AP}^q(Y)$ primitiv.

\textit{(AP1):} Sei $L = L_X \boxtimes 1 + 1 \boxtimes L_Y$ für ample Klassen $L_X \in \mathrm{Amp}(X)$, $L_Y \in \mathrm{Amp}(Y)$. Nach Lemma~\ref{lem:primitive_product} ist $\alpha \boxtimes \beta$ $L$-primitiv, also gleicht es seiner eigenen primitiven Komponente. Nach Lemma~\ref{lem:sign_product}:
\[
\begin{aligned}
&(-1)^{p+q} Q_L(\alpha \boxtimes \beta,\, \alpha \boxtimes \beta)\\
&\quad= \binom{N-2k}{n-2p}
   \bigl[(-1)^p Q_{L_X}(\alpha, \alpha)\bigr]\\
&\quad\cdot \bigl[(-1)^q Q_{L_Y}(\beta, \beta)\bigr].
\end{aligned}
\]
Der Binomialkoeffizient $\binom{N-2k}{n-2p} > 0$, und beide geklammerten Faktoren sind $\geq 0$ durch (AP1) für $\alpha$ und $\beta$. Daher gilt $(-1)^{p+q}\, Q_L(\alpha \boxtimes \beta, \alpha \boxtimes \beta) \geq 0$. Ample Klassen der Form $L_X \boxtimes 1 + 1 \boxtimes L_Y$ bilden einen Unterkegel von $\mathrm{Amp}(X \times Y)$. Da die Menge $\{L \in \overline{\mathrm{Amp}}(X \times Y) : (-1)^{p+q} Q_L(\gamma_0, \gamma_0) \geq 0\}$ abgeschlossen ist ($Q_L$ hängt stetig von $L$ ab) und diesen Unterkegel enthält, erstreckt sich die Ungleichung auf den Abschluss dieses Produktpolarisierungs-Unterkegels. Das genügt nur unter einer expliziten Néron-Severi-Vergleichshypothese, etwa wenn $N^1(X\times Y)_\mathbb{R}$ von den beiden Rückzugssummanden erzeugt wird und der Amplekegel entsprechend produktgeneriert ist. Picard-Zahl~$1$ eines einzelnen Faktors reicht im Allgemeinen nicht, da auf Produkten gemischte Divisorklassen auftreten können. Im Allgemeinen erfordert die Erweiterung auf ganz $\mathrm{Amp}(X \times Y)$ die Verifikation der Ungleichung für Polarisierungen mit gemischten Termen, was offen bleibt.

\textit{Konjugationskompatibilität:} Für beliebiges $\sigma \in \mathrm{Aut}(\mathbb{C})$ gilt $(\alpha \boxtimes \beta)^\sigma = \alpha^\sigma \boxtimes \beta^\sigma$ und $(X \times Y)^\sigma = X^\sigma \times Y^\sigma$. Daher lässt sich dieselbe (AP1)-Produktpolarisierungsrechnung nach Konjugation anwenden.

\textit{(AP3)-Vorbehalt:} Zwar gilt $\Lambda_L^j(\alpha \boxtimes \beta) = \sum_{a+b=j} \binom{j}{a}(\Lambda_{L_X}^a \alpha)\boxtimes(\Lambda_{L_Y}^b \beta)$, aber Positivität der Summanden kontrolliert nicht automatisch die primitive Projektion der Summe auf $X\times Y$. Die vollständige Kreuzterm-Analyse wird hier nicht durchgeführt. Deshalb ist die Proposition nur als Produktpolarisierungs-Vorzeichentest formuliert.
\end{proof}


\begin{remark}[Barriere-Lemma: HR-Positivität ist automatisch gesättigt]
\label{rem:barrier-lemma}
Bevor wir zum No-Go-Satz übergehen, heben wir den strukturellen Grund hervor, warum jede positivitätsbasierte Verstärkung der Hodge-Vermutung das Hodge-Riemann-Framework verlassen muss. Für eine glatte projektive Varietät $X$ und jede ample Klasse $L$ garantieren die Hodge-Riemann-Bilinearrelationen:
\begin{enumerate}
\item $(-1)^p Q_L(\alpha, \bar{\alpha}) > 0$ für jede nichtverschwindende primitive $(p,p)$-Klasse $\alpha$;
\item $Q_L$ ist nichtentartet auf $H^{p,p}_{\mathrm{prim}}(X)$.
\end{enumerate}
Diese Bedingungen sind \emph{automatisch erfüllt} für alle rationalen $(p,p)$-Klassen -- algebraische wie nichtalgebraische. Folglich kann keine Positivitätsbedingung, die durch die HR-Relationen impliziert wird, algebraische von nichtalgebraischen Klassen unterscheiden. Dies ist die Barriere: die ,,offensichtliche'' Verstärkung (Positivität unter allen Polarisierungen, allen Automorphismen, allen Lefschetz-Zerlegungen fordern) kollabiert zur Bedingung, eine absolute Hodge-Klasse zu sein (Satz~\ref{thm:nogo}). Die einzigen Auswege erfordern \emph{nichtbilineare} Positivitätsbedingungen (motivische Faltung, Galois-Mittelung, prismatische Einschränkungen) oder Informationen von \emph{außerhalb} der Hodge-Riemann-Struktur (o-Minimalität, derivierte Kategorien).
\end{remark}

% ===================================================================
% 6. DAS NO-GO-THEOREM
% ===================================================================
\section{Der No-Go-Satz}
\label{sec:nogo}

Wir formulieren nun das zentrale Resultat: die Hodge-Vermutung, umformuliert als Positivitätsobstruktion.

\begin{theorem}[Positivitätsobstruktion -- Bedingt]
\label{thm:nogo}
Es gelte die Arithmetische Positivitätsvermutung~\ref{conj:AP}. Sei $X$ eine glatte projektive Varietät über $\mathbb{C}$, und sei $\alpha \in H^{2p}(X, \mathbb{Q}) \cap H^{p,p}(X)$ eine rationale Hodge-Klasse, die \textbf{keine} $\mathbb{Q}$-Linearkombination algebraischer Zykel ist. Dann existiert mindestens eine der folgenden Obstruktionen:

\begin{enumerate}
\item[\textbf{(O1)}] \textbf{Polarisierungsinstabilität:} Es existiert eine ample Klasse $L \in \mathrm{Amp}(X)$, sodass die primitive Komponente $\alpha_0$ von $\alpha$ die Bedingung $(-1)^p\,Q_L(\alpha_0, \alpha_0) < 0$ erfüllt.

\item[\textbf{(O2)}] \textbf{Galois-Instabilität:} Es existiert $\sigma \in \mathrm{Aut}(\mathbb{C})$, sodass entweder $\alpha^\sigma$ keine Hodge-Klasse auf $X^\sigma$ ist, oder $\alpha^\sigma$ (AP1) auf $X^\sigma$ verletzt.

\item[\textbf{(O3)}] \textbf{Lefschetz-Instabilität:} Es existieren ein zulässiges $j \geq 1$ im Bereich von (AP3) und $L \in \mathrm{Amp}(X)$, sodass $(-1)^{p+j}\,Q_L(\gamma_{j,0}, \gamma_{j,0}) < 0$, wobei $\gamma_{j,0}$ die $L$-primitive Komponente von $\Lambda_L^j \alpha \in H^{2(p+j)}(X)$ ist.
\end{enumerate}
\end{theorem}

\begin{proof}
Durch Kontraposition: wenn $\alpha$ (AP1)--(AP3) erfüllt, dann gilt $\alpha \in \mathrm{AP}^p(X) = \mathrm{cl}(\mathrm{CH}^p(X)_\mathbb{Q})$ durch die APV, im Widerspruch zur Annahme, dass $\alpha$ nicht algebraisch ist.
\end{proof}

\begin{remark}[Struktur der Obstruktion]
\label{rem:obstruction_structure}
Die drei Obstruktionen müssen im Licht der Hodge-Riemann-Relationen interpretiert werden.

\textbf{(O1) und (O3) sind für $(p,p)$-Klassen vakuos.} Die Hodge-Riemann-Bilinearrelationen (Satz~\ref{thm:HR}) garantieren $(-1)^p Q_L(\alpha_0, \alpha_0) \geq 0$ für \textit{jede} rationale $(p,p)$-Klasse $\alpha$, nicht nur für algebraische. Die primitive Komponente $\alpha_0$ liegt in $P^{p,p}(X)$, und die Hodge-Riemann-Form ist dort positiv definit (mit dem Vorzeichen $(-1)^p$). Analog liegt für zulässige AP3-Potenzen die primitive Komponente $\gamma_{j,0}$ von $\Lambda_L^j \alpha$ in $P^{p+j,p+j}(X)$, sodass $(-1)^{p+j} Q_L(\gamma_{j,0}, \gamma_{j,0}) \geq 0$ automatisch gilt. Daher \textit{können Obstruktionen (O1) und (O3) für rationale Hodge-Klassen im angegebenen AP-Bereich niemals eintreten}.

\textbf{(O2) ist die einzige effektive Obstruktion.} Der einzige Weg, auf dem eine Hodge-Klasse arithmetische Positivität verletzen kann, ist über (O2a): Versagen der Absolutheit ($\alpha^\sigma$ ist keine Hodge-Klasse auf $X^\sigma$ für ein $\sigma$). Teil (O2b) -- Versagen der Hodge-Riemann-Positivität auf $X^\sigma$ -- ist ebenfalls vakuos, da $\alpha^\sigma \in H^{p,p}(X^\sigma)$ die Hodge-Riemann-Ungleichungen automatisch erfüllt.

Folglich reduziert sich der No-Go-Satz auf Delignes Beobachtung: nichtabsolute Hodge-Klassen (falls sie existieren) können nicht algebraisch sein.
\end{remark}

\subsection{Die unbedingte Aussage}

Während Satz~\ref{thm:nogo} von der APV abhängt, können wir ein unbedingtes Teilresultat formulieren:

\begin{proposition}[Unbedingte Obstruktion]
\label{prop:unconditional}
Sei $\alpha \in H^{2p}(X, \mathbb{Q}) \cap H^{p,p}(X)$ eine rationale Hodge-Klasse. Wenn $\alpha$ (AP2a) nicht erfüllt -- d.h., es existiert $\sigma \in \mathrm{Aut}(\mathbb{C})$ mit $\alpha^\sigma \notin H^{p,p}(X^\sigma)$ -- dann ist $\alpha$ nicht algebraisch.
\end{proposition}

\begin{proof}
Dies ist Delignes Resultat (Satz~\ref{thm:deligne_absolute}, Teil~1, Kontraposition): algebraische Klassen sind absolut.
\end{proof}

Wie in Abschnitt~\ref{sec:discussion} (Bemerkung~\ref{rem:ap_equals_absolute}) gezeigt, werden die Hodge-Riemann-Positivitätsbedingungen (AP1), (AP2b) und (AP3) von jeder absoluten Hodge-Klasse automatisch erfüllt. Der Inhalt der APV jenseits von Proposition~\ref{prop:unconditional} ist daher identisch mit Delignes Frage: \textit{Sind alle absoluten Hodge-Klassen algebraisch?} Die ,,Lücke'' zwischen absoluten Hodge-Klassen und algebraischen Klassen wird durch die Positivitätsbedingungen dieser Arbeit nicht verkleinert.


% ===================================================================
% 7. DER FALL ABELSCHER VARIETAETEN
% ===================================================================
\section{Abelsche Varietäten}
\label{sec:abelian}

Abelsche Varietäten bieten die günstigste Umgebung für den arithmetischen Positivitätsansatz, weil Delignes Satz (alle Hodge-Klassen sind absolut) die Obstruktion (O2a) eliminiert.

\begin{proposition}[Äquivalenz für abelsche Varietäten]
\label{thm:abelian}
Für abelsche Varietäten ist die Arithmetische Positivitätsvermutung~\ref{conj:AP} äquivalent zur Hodge-Vermutung.
\end{proposition}

\begin{proof}
Sei $A$ eine abelsche Varietät der Dimension $g$, und sei $\alpha \in H^{2p}(A, \mathbb{Q}) \cap H^{p,p}(A)$.

Nach Delignes Satz (Satz~\ref{thm:deligne_absolute}, Teil~2) ist jede Hodge-Klasse auf einer abelschen Varietät absolut. Proposition~\ref{prop:absolute_implies_ap} liefert daher
\[
H^{2p}(A,\mathbb{Q})\cap H^{p,p}(A)=\mathrm{AbsHodge}^p(A)=\mathrm{AP}^p(A).
\]
Damit ist die APV für $A$,
\[
\mathrm{AP}^p(A)=\mathrm{cl}(\mathrm{CH}^p(A)_{\mathbb{Q}}),
\]
wortgleich die Hodge-Vermutung für $A$,
\[
H^{2p}(A,\mathbb{Q})\cap H^{p,p}(A)
=\mathrm{cl}(\mathrm{CH}^p(A)_{\mathbb{Q}}).
\]
Dies beweist die Äquivalenz. Es beweist nicht die Hodge-Vermutung für allgemeine abelsche Varietäten; es zeigt nur, dass arithmetische Positivität keine zusätzliche Bedingung hinzufügt, sobald Delignes Absolutheitssatz verfügbar ist.
\end{proof}

\begin{remark}[Bekannte Positivfälle]
Für abelsche Varietäten, bei denen alle Hodge-Klassen als algebraisch bekannt sind (etwa mehrere niedrigdimensionale, CM- oder Produktfälle in der Literatur), ist damit auch die äquivalente APV-Aussage verifiziert. Allgemein liefert Delignes Satz jedoch Absolutheit, nicht Algebraizität; offen bleibt zu zeigen, dass absolute Hodge-Klassen durch algebraische Zykel repräsentiert werden.
\end{remark}


% ===================================================================
% 8. DIE PRAEZISE OBSTRUKTION
% ===================================================================
\section{Die präzise Obstruktion}
\label{sec:obstruction}

Wir identifizieren nun den exakten Mechanismus, durch den nichtalgebraische Hodge-Klassen arithmetische Positivität verletzen.

\subsection{Die Galois-Hodge-Riemann-Form}

\begin{definition}[Galois-Hodge-Riemann-Form]
\label{def:galois_HR}
Für $\sigma \in \mathrm{Aut}(\mathbb{C})$ und $L \in \mathrm{Amp}(X)$ definiere:
\begin{equation}
Q_L^\sigma(\alpha) := Q_{L^\sigma}(\alpha^\sigma, \alpha^\sigma)\,.
\label{eq:galois_HR}
\end{equation}
Das \textit{Galois-Hodge-Riemann-Spektrum} von $\alpha$ ist:
\begin{equation}
\begin{aligned}
\mathrm{Spec}_{\mathrm{GHR}}(\alpha)
:= \{\, &(-1)^p Q_L^\sigma(\alpha) :\\
&\sigma \in \mathrm{Aut}(\mathbb{C}),\,
  L \in \mathrm{Amp}(X)\,\} \subset \mathbb{R}.
\end{aligned}
\end{equation}
\end{definition}

\begin{proposition}[Algebraische Klassen haben nichtnegatives GHR-Spektrum]
\label{prop:ghr_positive}
Wenn $\alpha = \mathrm{cl}(Z)$ für einen algebraischen Zykel $Z$, dann gilt:
\begin{equation}
\mathrm{Spec}_{\mathrm{GHR}}(\alpha) \subseteq [0, \infty)\,.
\end{equation}
\end{proposition}

\begin{proof}
Algebraische Zykel transformieren kovariant unter $\sigma$: $\mathrm{cl}(Z)^\sigma = \mathrm{cl}(Z^\sigma)$, und $Z^\sigma$ ist ein algebraischer Zykel auf $X^\sigma$. Die primitive Komponente von $\mathrm{cl}(Z^\sigma)$ liegt in $P^{p,p}(X^\sigma)$, sodass die Hodge-Riemann-Relationen $(-1)^p\, Q_{L^\sigma}(\mathrm{cl}(Z^\sigma)_0, \mathrm{cl}(Z^\sigma)_0) \geq 0$ liefern, was die geforderte Nichtnegativität des GHR-Spektrums ergibt.
\end{proof}

\begin{conjecture}[GHR-Obstruktionsvermutung]
\label{conj:ghr}
Wenn $\alpha$ eine rationale Hodge-Klasse ist, die \textbf{nicht} algebraisch ist, dann gilt:
\begin{equation}
\mathrm{Spec}_{\mathrm{GHR}}(\alpha) \cap (-\infty, 0) \neq \emptyset\,.
\end{equation}
Das heißt, es existieren $\sigma$ und $L$, sodass $Q_L^\sigma(\alpha) < 0$.
\end{conjecture}

Dies ist eine präzisere Version der Arithmetischen Positivitätsvermutung: sie identifiziert die Obstruktion als das Versagen der Galois-Hodge-Riemann-Form, nichtnegativ zu bleiben.

\begin{remark}[Logischer Status der GHR-Obstruktionsvermutung]
\label{rem:ghr_status}
Die GHR-Obstruktionsvermutung muss im Licht von $\mathrm{AP}^p = \mathrm{AbsHodge}^p$ (Bemerkung~\ref{rem:ap_equals_absolute}) vorsichtig interpretiert werden. Wenn $\alpha$ eine absolute Hodge-Klasse ist, dann gilt $(-1)^p Q_L^\sigma(\alpha) \geq 0$ für alle $\sigma$ und $L$ nach Proposition~\ref{prop:absolute_implies_ap}. Daher kann die Vermutung nur für nichtalgebraische Hodge-Klassen ,,feuern'', die \textit{nicht absolut} sind -- in diesem Fall wird die Obstruktion bereits durch (AP2a) (Versagen der Absolutheit) bezeugt, nicht durch Negativität von $Q_L^\sigma$. Falls nichtalgebraische absolute Hodge-Klassen existieren (d.h., falls Delignes Frage eine negative Antwort hat), wäre die GHR-Obstruktionsvermutung für solche Klassen \textit{falsch}, da ihr GHR-Spektrum automatisch nichtnegativ ist. Die Vermutung ist daher äquivalent zur Aussage, dass nichtalgebraische absolute Hodge-Klassen nicht existieren -- was präzise Delignes Frage ist.
\end{remark}

\begin{remark}[Falsifizierbarkeit der GHR-Diagnostik]
\label{rem:ghr-falsifiability}
Das GHR-Spektrum kann als diagnostisches Werkzeug mit klaren Falsifikationskriterien verwendet werden:
\begin{enumerate}
\item \emph{Positiver Test:} Falls eine Klasse $\alpha$ die Eigenschaft $(-1)^p Q_L^\sigma(\alpha^\sigma) < 0$ für ein $\sigma \in \mathrm{Aut}(\mathbb{C})$ und ein $L$ besitzt, dann ist $\alpha$ beweisbar nichtalgebraisch (nach Satz~\ref{thm:easy_direction}).
\item \emph{Null-Test:} Falls $(-1)^p Q_L^\sigma(\alpha^\sigma) \geq 0$ für alle $\sigma$ und alle $L$, dann ist $\alpha$ arithmetisch positiv -- somit absolut Hodge -- aber nicht notwendigerweise algebraisch. Dies ist der Bereich, in dem die GHR-Diagnostik \emph{stumm} ist.
\item \emph{Struktureller Test:} Das GHR-Spektrum detektiert Nichtalgebraizität nur für Klassen, die \emph{nicht} absolut Hodge sind. Für absolute nichtalgebraische Klassen (falls sie existieren) wird eine strikt stärkere Diagnostik benötigt.
\end{enumerate}
Das GHR-Spektrum fungiert somit als \emph{notwendiger-Bedingung-Filter}: es kann Nichtalgebraizität zertifizieren, aber nicht Algebraizität. Diese Asymmetrie spiegelt die Pattern-A-Architektur im gesamten FST-Programm wider, wo Positivitätszertifikate Alternativen ausschließen, aber keine Lösungen konstruieren.
\end{remark}

\begin{example}[GHR-Spektrum der Diagonalklasse]
\label{ex:diagonal}
Sei $X = \mathbb{P}^n$ und betrachte die Diagonaleinbettung $\Delta: \mathbb{P}^n \hookrightarrow \mathbb{P}^n \times \mathbb{P}^n$. Die Klasse $[\Delta] \in H^{2n}(\mathbb{P}^n \times \mathbb{P}^n, \mathbb{Q})$ ist algebraisch (sie ist die Zykelklasse der Diagonalen). Nach Satz~\ref{thm:easy_direction} ist $[\Delta]$ arithmetisch positiv. Ihre Künneth-Zerlegung ist $[\Delta] = \sum_{k=0}^n h^k \boxtimes h^{n-k}$, wobei $h$ die Hyperebenenklasse ist. Da $\mathbb{P}^n$ keine nichttrivialen Automorphismen von $\mathbb{C}$ hat, die auf seiner Kohomologie wirken (die Hodge-Struktur ist rein Tate), fixiert jedes $\sigma \in \mathrm{Aut}(\mathbb{C})$ die Klasse $[\Delta]^\sigma = [\Delta]$. Die Selbstschnittzahl ist $\int_{X \times X} [\Delta]^2 = \chi(\mathbb{P}^n) = n + 1$ (die Euler-Charakteristik). Da $[\Delta]$ algebraisch ist, erfüllt ihre primitive Komponente $[\Delta]_0$ die Ungleichung $(-1)^n Q_L([\Delta]_0, [\Delta]_0) \geq 0$ nach den Hodge-Riemann-Relationen. Wir bemerken, dass $[\Delta]$ \emph{nicht} primitiv ist: $L \cup [\Delta] \neq 0$ im Allgemeinen, also ist die primitive Zerlegung nichttrivial. Dieses Beispiel ist ,,trivial positiv'', weil es algebraisch ist, aber es illustriert die Berechnung: für eine hypothetische nichtalgebraische Hodge-Klasse auf einer komplexeren Varietät würde man $\sigma$ und $L$ suchen, bei denen die GHR-Form negativ wird.
\end{example}

\begin{example}[Nichttriviales GHR-Spektrum: CM-Endomorphismenklasse]
\label{ex:cm_endomorphism}
Betrachte die elliptische Kurve $E: y^2 = x^3 - x$ über $\mathbb{Q}(i)$, die komplexe Multiplikation durch $\mathbb{Z}[i]$ besitzt, mit $\mathrm{End}(E) \otimes \mathbb{Q} = \mathbb{Q}(i)$. Der CM-Endomorphismus $\tau = i$ (Multiplikation mit $i$) hat einen Graphen $\Gamma_\tau \subset E \times E$, der eine algebraische Klasse $[\Gamma_\tau] \in H^2(E \times E, \mathbb{Q}) \cap H^{1,1}(E \times E)$ liefert. Diese Klasse liegt \emph{außerhalb} des Lefschetz-Rings: der Lefschetz-Ring von $E \times E$ (erzeugt durch die beiden Faserklassen $[E \times \{0\}]$, $[\{0\} \times E]$ und die Diagonale $[\Delta]$) enthält $[\Gamma_\tau]$ nicht, da $\tau$ nichttrivial auf $H^{1,0}(E)$ wirkt.

Das GHR-Spektrum von $[\Gamma_\tau]$ ist nichttrivial. Sei $\sigma \in \mathrm{Aut}(\mathbb{C})$ die auf $\mathbb{Q}(i)$ eingeschränkte komplexe Konjugation, also $\sigma(i) = -i$. Dann gilt $E^\sigma \cong E$ (da $y^2 = x^3 - x$ rationale Koeffizienten hat), aber der CM-Endomorphismus transformiert sich als $\tau^\sigma = \sigma(i) = -i$, also $[\Gamma_\tau]^\sigma = [\Gamma_{-\tau}]$. Sei $L = L_1 \boxtimes 1 + 1 \boxtimes L_1$ die Produktpolarisierung, wobei $L_1$ die Hauptpolarisierung auf $E$ ist. Dann:
\begin{align}
(-1)^1 Q_L^\mathrm{id}([\Gamma_\tau])
&= (-1)^1 Q_L([\Gamma_i], [\Gamma_i]) \notag\\
&= 1 + |\tau|^2 = 2, \label{eq:cm_id}\\
(-1)^1 Q_L^\sigma([\Gamma_\tau])
&= (-1)^1 Q_{L^\sigma}([\Gamma_{-i}], [\Gamma_{-i}]) \notag\\
&= 1 + |-\tau|^2 = 2. \label{eq:cm_sigma}
\end{align}
Beide Werte sind positiv, wie von Satz~\ref{thm:easy_direction} garantiert. Allerdings unterscheidet sich die \emph{Zwischenberechnung}: die Klassen $[\Gamma_i]$ und $[\Gamma_{-i}]$ haben verschiedene Künneth-Zerlegungen in $H^1(E) \otimes H^1(E)$, was die Tatsache widerspiegelt, dass $\sigma$ die Eigenräume der CM-Aktion auf $H^1(E, \mathbb{C}) = H^{1,0}(E) \oplus H^{0,1}(E)$ permutiert. Explizit wirkt $\tau^*$ auf $H^{1,0}(E)$ durch Multiplikation mit $i$ und auf $H^{0,1}(E)$ durch $-i$; unter $\sigma$ werden diese Eigenwerte vertauscht.

Der wesentliche Punkt ist nicht eine Änderung der beiden angezeigten Zahlenwerte: In diesem symmetrischen CM-Beispiel sind beide gleich~$2$. Nichttrivial ist, dass sich die Zwischenzerlegung nach Künneth unter $\sigma$ ändert, während der vorzeichenbehaftete Hodge-Riemann-Wert nichtnegativ bleibt. Das Beispiel ist daher ein Konsistenzcheck für die GHR-Visualisierung, kein Algebraizitätskriterium: Die Algebraizität von $[\Gamma_\tau]$ stammt aus dem CM-Graphen, während arithmetische Positivität automatisch ist, sobald die Klasse absolut ist.
\end{example}

\subsection{Warum die Obstruktion natürlich ist}

Die GHR-Obstruktion ist aus folgendem Grund natürlich. Die Hodge-Riemann-Form $Q_L$ hängt von beidem ab:
\begin{itemize}
\item der \textit{komplexen Struktur} von $X$ (durch die Hodge-Zerlegung), und
\item der \textit{algebraischen Struktur} (durch die Polarisierung $L$).
\end{itemize}

Für algebraische Klassen sind die komplexe und die algebraische Struktur ,,ausgerichtet'': die Klasse stammt von einer Untervarietät, die gleichzeitig eine komplexe Untermannigfaltigkeit und ein algebraischer Zykel ist. Die Galois-Aktion $\sigma$ kann die komplexe und algebraische Struktur ,,fehlausrichten'', aber für algebraische Klassen ,,richtet'' der Zykel $Z^\sigma$ sie auf $X^\sigma$ wieder aus.

Für eine \textit{nichtalgebraische} Hodge-Klasse gibt es keinen Zykel, der diese Neuausrichtung durchführt. Die Klasse existiert in der komplexen Struktur von $X$, hat aber keinen algebraischen Anker. Wenn $\sigma$ die komplexe Struktur verdrillt, kann die Positivität der Hodge-Riemann-Form versagen, weil es nichts gibt, das die Verdrillung kompensiert.

Dies ist der präzise Mechanismus der Obstruktion: \textit{algebraische Klassen sind starr genug, um alle Galois-Verdrillungen mit intakter Positivität zu überstehen, während transzendente Klassen es nicht sind.} Wir betonen jedoch, dass nach Bemerkung~\ref{rem:ap_equals_absolute} und Bemerkung~\ref{rem:obstruction_structure} die Hodge-Riemann-Positivität hier keine eigenständige Rolle spielt: die einzige effektive Obstruktion ist Versagen der Absolutheit (O2a). Die ,,Ausrichtungs''-Erzählung beschreibt den Mechanismus der Absolutheit, nicht der Hodge-Riemann-Positivität.


% ===================================================================
% 9. DISKUSSION
% ===================================================================
\section{Diskussion}
\label{sec:discussion}

\subsection{Die verbleibende Lücke}

Die Arbeit etabliert:
\begin{enumerate}
\item Algebraische Klassen $\Rightarrow$ arithmetisch positiv (Satz~\ref{thm:easy_direction}, unbedingt).
\item $\mathrm{AP}^p(X) = \mathrm{AbsHodge}^p(X)$ (Bemerkung~\ref{rem:ap_equals_absolute}, unbedingt).
\item Begrenzte Funktorialitätstests für Rückzugsunterkegel endlicher Morphismen und primitive Künneth-Vorzeichenformeln unter Produktpolarisierungen; volle Produkt-AP-Mitgliedschaft folgt nicht aus der direkten Vorzeichenrechnung allein (Abschnitt~\ref{sec:functoriality}).
\item Identifikation des Galois-Hodge-Riemann-Spektrums als diagnostisches Werkzeug (Abschnitt~\ref{sec:obstruction}).
\end{enumerate}

Als Konsequenz von Punkt~2 ist die verbleibende Lücke präzise Delignes Frage:

\begin{quote}
\textit{Sind alle absoluten Hodge-Klassen algebraisch?}
\end{quote}

Dies ist ein langstehendes offenes Problem. Das Positivitäts-Framework löst es nicht, bietet aber einen neuen Blickwinkel: die No-Go-Formulierung und das GHR-Spektrum machen explizit, welche Galois-Operationen ein Versagen der Absolutheit bezeugen könnten. Sie unterscheiden algebraische Klassen nicht von nichtalgebraischen absoluten Hodge-Klassen, falls solche Klassen existieren.

\subsection{Beziehung zu Delignes Programm}

Delignes Framework absoluter Hodge-Klassen \cite{Deligne1982} liefert Bedingung (AP2a). Die vorliegende Arbeit beabsichtigte ursprünglich, Delignes Absolutheit durch Hinzufügen zu verstärken:

\begin{enumerate}
\item \textbf{Hodge-Riemann-Positivität unter Konjugation} (AP2b): Nicht nur der Hodge-\textit{Typ}, sondern auch das Hodge-Riemann-\textit{Vorzeichen} muss erhalten bleiben.

\item \textbf{Universelle Polarisierung} (AP1 für alle $L$): Alle amplen Klassen müssen nichtnegative Hodge-Riemann-Form ergeben.
\end{enumerate}

Jedoch zeigen Proposition~\ref{prop:absolute_implies_ap} und Bemerkung~\ref{rem:ap_equals_absolute}, dass diese Bedingungen für jede absolute Hodge-Klasse \textit{automatisch} durch die Hodge-Riemann-Bilinearrelationen erfüllt sind. Die Bedingungen (AP1), (AP2b) und (AP3) schränken daher nicht über Delignes Absolutheit hinaus ein. Die Frage ,,Sind alle absoluten Hodge-Klassen algebraisch?'' bleibt präzise Delignes motivierendes Problem, und die APV ist äquivalent dazu.

\begin{proposition}
\label{prop:absolute_implies_ap}
Auf einer glatten projektiven Varietät $X/\mathbb{C}$ ist jede absolute Hodge-Klasse arithmetisch positiv:
$\mathrm{AbsHodge}^p(X) \subseteq \mathrm{AP}^p(X)$.
\end{proposition}

\begin{proof}
Sei $\alpha$ eine absolute Hodge-Klasse. (AP1): Die primitive Komponente $\alpha_0$ bezüglich einer amplen Klasse $L$ liegt in $P^{p,p}(X) \cap H^{2p}(X, \mathbb{Q})$. Für rationale $(p,p)$-Klassen gilt $\bar{\alpha}_0 = \alpha_0$ und $i^{p-p} = 1$, sodass $h_L(\alpha_0, \alpha_0) = (-1)^{p(2p-1)} Q_L(\alpha_0, \alpha_0) = (-1)^p Q_L(\alpha_0, \alpha_0)$. Die HR-Relationen geben $h_L > 0$, also $(-1)^p Q_L(\alpha_0, \alpha_0) > 0$ für $\alpha_0 \neq 0$. (AP2a) gilt nach Voraussetzung. (AP2b) folgt, da $\alpha^\sigma \in H^{p,p}(X^\sigma)$ für absolutes $\alpha$, und das (AP1)-Argument auf $X^\sigma$ anwendbar ist. (AP3) folgt für zulässige $j$ im Bereich von Definition~\ref{def:arith_pos} analog, da die primitive Komponente von $\Lambda_L^j \alpha$ in $P^{p+j,p+j}(X)$ liegt.
\end{proof}

\begin{remark}[Kritische Beobachtung]
\label{rem:ap_equals_absolute}
Proposition~\ref{prop:absolute_implies_ap} zeigt $\mathrm{AbsHodge}^p(X) \subseteq \mathrm{AP}^p(X)$. Die umgekehrte Inklusion $\mathrm{AP}^p(X) \subseteq \mathrm{AbsHodge}^p(X)$ gilt nach Definition, da Bedingung (AP2a) Absolutheit fordert. Daher:
\begin{equation}
\mathrm{AP}^p(X) = \mathrm{AbsHodge}^p(X)\,.
\end{equation}
Die Bedingungen (AP1), (AP2b) und (AP3) stellen \textit{keine} zusätzliche Einschränkung jenseits von Delignes Absolutheit dar. Die Hodge-Riemann-Bilinearrelationen \textit{garantieren} die Positivitätsbedingungen für jede Klasse vom Typ $(p,p)$ automatisch.

Dies hat eine tiefgreifende Konsequenz für die APV: die Arithmetische Positivitätsvermutung $\mathrm{AP}^p(X) = \mathrm{cl}(\mathrm{CH}^p(X)_\mathbb{Q})$ ist \textbf{äquivalent} zur Aussage, dass jede absolute Hodge-Klasse algebraisch ist -- was präzise Delignes ursprüngliche Frage \cite{Deligne1982} ist. Das in dieser Arbeit eingeführte Positivitäts-Framework liefert daher keine genuin neue Vermutung, sondern vielmehr eine neue \textit{Perspektive} auf Delignes Frage durch die Sprache der Hodge-Riemann-Positivität und No-Go-Obstruktionen.

Die GHR-Obstruktion (Definition~\ref{def:galois_HR}) bleibt als diagnostisches Werkzeug sinnvoll: sie macht explizit, \textit{welche} Galois-Konjugation für eine nichtalgebraische absolute Hodge-Klasse versagen müsste. Der mathematische Kerninhalt reduziert sich jedoch auf: \textit{Sind alle absoluten Hodge-Klassen algebraisch?}
\end{remark}

Die unbedingte Kette lautet:
\begin{equation}
\begin{aligned}
\text{Algebraisch}
&\xRightarrow{\text{Thm.~\ref{thm:easy_direction}}}
  \text{AP} = \text{Absolut}\\
&\xhookrightarrow{\text{Definition}} \text{Hodge}\,,
\end{aligned}
\end{equation}
wobei die Gleichheit $\text{AP} = \text{Absolut}$ durch Bemerkung~\ref{rem:ap_equals_absolute} gegeben ist. Die HV verlangt die zusammengesetzte Rückinklusion $\text{Hodge}\subseteq\text{Algebraisch}$; diese zerfällt in zwei logisch verschiedene Fragen: ob alle Hodge-Klassen absolut sind und ob alle absoluten Hodge-Klassen algebraisch sind. Für abelsche Varietäten liefert Delignes Satz die erste Rückinklusion, die zweite bleibt das eigentliche Algebraizitätsproblem.

\subsection{Verbindung zum breiteren Programm}

Dieser Ansatz fügt sich in ein breiteres Muster ein, das über fundamentale Mathematik und Physik beobachtet wird:

\begin{center}
\resizebox{\columnwidth}{!}{%
\begin{tabular}{lll}
\toprule
Problem & Positivitätskonzept & Resultat \\
\midrule
Weil-Verm. & Frobenius-Eigenwert & Deligne 1974~\cite{Deligne1974} \\
BSD (Rang 1) & N\'eron-Tate-Höhe & Gross-Zagier 1986~\cite{GrossZagier1986} \\
\textbf{Hodge} & \textbf{GHR-Spektrum} & \textbf{Diese Arbeit (bedingt)} \\
\bottomrule
\end{tabular}}
\end{center}

In jedem Fall besteht die Strategie nicht darin, das gewünschte Objekt zu konstruieren, sondern zu zeigen, dass seine Nichtexistenz eine Positivitätsbedingung verletzen würde. Die vorliegende Arbeit wendet diese Strategie auf die Hodge-Vermutung an. Frühere Arbeiten des Autors \cite{GeigerRH2026c, FST-Unified2026} erforschten ähnliche positivitätsbasierte Ansätze in anderen Kontexten.

\begin{remark}[Meta-Muster: versuchte Verstärkung kollabiert an der Sättigung]
\label{rem:meta-pattern-saturation}
Die strukturelle Lektion dieser Arbeit geht über die Hodge-Vermutung hinaus. In jeder Pattern-A-Instanziierung im FST-Programm ist die ,,offensichtliche'' Verstärkung der Positivitätsbedingung automatisch durch das Strukturtheorem gesättigt:
\begin{itemize}
\item \emph{Hodge:} HR-Positivität $\Rightarrow$ AP $=$ AbsHodge (diese Arbeit).
\item \emph{NS-Skelett:} Lipschitz-Regularität der Attraktor-Projektion $\Rightarrow$ erfordert positive Reichweite $\Rightarrow$ äquivalent zur Existenz einer Inertialmannigfaltigkeit (No-Go-Zone, vgl.\ die Begleitarbeit zu NS).
\item \emph{BSD:} Höhenpositivität für Rang $\leq 1$ $\Rightarrow$ gesättigt durch Gross--Zagier + Kolyvagin; die ,,nächste'' Positivität (Rang $\geq 2$) erfordert Objekte, die noch nicht existieren.
\end{itemize}
In jedem Fall erfordert Fortschritt das \emph{Verlassen} des gesättigten Positivitäts-Frameworks -- entweder durch Einführung nichtlinearer/nichtbilinearer Bedingungen oder durch Import von Regularität aus einer anderen Quelle (z.B.\ die TLL/LDI-Maschinerie für NS oder prismatische Kohomologie für Hodge). Das Sättigungsphänomen ist kein Versagen der Methode, sondern eine \emph{Diagnostik}: es identifiziert die präzise Grenze, an der die aktuellen Werkzeuge ihre Reichweite erschöpfen.
\end{remark}

\subsection{Einschränkungen}

\begin{enumerate}
\item \textbf{Die schwere Richtung ist offen.} Diese Arbeit beweist nicht die Hodge-Vermutung. Sie formuliert sie als Positivitätsaussage um und identifiziert die präzise Obstruktion, aber der Beweis, dass arithmetische Positivität Algebraizität erzwingt, bleibt die zentrale Herausforderung. Ferner impliziert die APV ($\mathrm{AP}^p = \mathrm{alg}^p$) allein nicht die HV ($\mathrm{Hodge}^p = \mathrm{alg}^p$): man benötigt zusätzlich, dass alle Hodge-Klassen arithmetisch positiv sind (siehe Proposition~\ref{prop:equivalence}).

\item \textbf{Teilweise Skizzen verbleiben in Abschnitt~\ref{sec:functoriality}.} Die Künneth-Rechnung ist ein primitiver Produktpolarisierungs-Vorzeichentest (Proposition~\ref{prop:kunneth}, Lemmas~\ref{lem:primitive_product}--\ref{lem:sign_product}), kein eigenständiger Beweis voller AP-Funktorialität. AP3 für Lefschetz-Potenzen auf $X\times Y$, die Erweiterung auf nichtprimitive Klassen und der Vorwärtsschub für gefaserte Morphismen erfordern sorgfältige Behandlung der Kreuzterme in der primitiven Zerlegung.

\item \textbf{Der AP-Kegel gleicht dem absoluten Hodge-Kegel.} Wie in Bemerkung~\ref{rem:ap_equals_absolute} gezeigt, gilt $\mathrm{AP}^p(X) = \mathrm{AbsHodge}^p(X)$, sodass die APV auf Delignes Frage reduziert. Das Positivitäts-Framework liefert keinen neuen mathematischen Inhalt jenseits einer neuen Sprache und diagnostischer Werkzeuge (das GHR-Spektrum) für das bestehende Problem.

\item \textbf{Verbleibende Lücke bei abelschen Varietäten.} Für abelsche Varietäten identifiziert Delignes Satz alle Hodge-Klassen mit absoluten Hodge-Klassen; AP fällt also mit dem ganzen Hodge-Klassenraum zusammen. Offen bleibt nicht Absolutheit, sondern Algebraizität dieser absoluten Klassen; dies ist nur in Spezialfamilien bekannt.
\end{enumerate}

\subsection{Ausblick}

Vier Richtungen sind unmittelbar:

\begin{enumerate}
\item \textbf{Verifikation der GHR-Obstruktionsvermutung für bekannte nichtalgebraische Klassen.} Ein Testfall ist Voisins Gegenbeispiel zur Hodge-Vermutung für K\"ahler-Mannigfaltigkeiten \cite{Voisin2002b}. Wie in Bemerkung~\ref{rem:ghr_status} festgestellt, kann die Vermutung allerdings nur Nichtabsolutheit detektieren; für absolute Hodge-Klassen, die möglicherweise nicht algebraisch sind, ist das GHR-Spektrum automatisch nichtnegativ.

\item \textbf{Verstärkung der AP-Bedingungen.} Da $\mathrm{AP}^p = \mathrm{AbsHodge}^p$ (Bemerkung~\ref{rem:ap_equals_absolute}), geht die aktuelle Definition arithmetischer Positivität nicht über Delignes Absolutheit hinaus. Man könnte \textit{zusätzliche} Positivitätsbedingungen suchen -- jenseits der durch die Hodge-Riemann-Relationen implizierten -- die die Lücke zwischen absoluten Hodge-Klassen und algebraischen Klassen einengen könnten. Kandidaten umfassen Bedingungen aus höheren Hodge-Riemann-Formen, $p$-adischer Hodge-Theorie oder motivischen Beschränkungen.

\item \textbf{Vergleiche die Kategorie der arithmetisch positiven Klassen} mit bestehenden Formalismen absoluter Hodge- und motivierter Klassen. Da $\mathrm{AP}=\mathrm{AbsHodge}$ gilt, ist sie kein Substitut für Motive; ihre Rolle ist diagnostisch und vergleichend, insbesondere für funktorielle Operationen und Kompatibilität mit den Standardvermutungen.

\item \textbf{Der diagnostische quadratische Lokus.} Der Hodge-Riemann-Teil der AP-Definition ist durch \emph{quadratische} Ungleichungen ausgedrückt (die Hodge-Riemann-Form ist bilinear, also sind die Positivitätsbedingungen quadratisch in $\alpha$). Da der volle AP-Lokus nach Bemerkung~\ref{rem:ap_equals_absolute} gleich $\mathrm{AbsHodge}^p(X)$ ist, ist dieser quadratische Lokus kein separates Algebraizitätskriterium. Er ist ein diagnostischer Schatten der Positivitätsbedingungen: abgeschlossen unter positiver Skalierung und symmetrisch, aber seine Randstruktur zeigt nur, wo die Hodge-Riemann-Visualisierung degeneriert. Sie klärt die Grenzen HR-basierter Diagnostik, nicht die schwere Richtung der APV.
\end{enumerate}


\begin{remark}[Arakelov-Brücke zu BSD]
\label{rem:arakelov-bridge}
Die Hodge--Riemann-Form $(-1)^p Q_L$ lässt sich als Positivität an der archimedischen (unendlichen) Primstelle in der Arakelov-Schnittpaarung interpretieren. Zusammen mit der N\'eron--Tate-Höhe (Positivität an endlichen Primstellen, vgl.\ das Begleitpaper zu BSD) legt dies ein vereinheitlichtes arithmetisches Positivitätsprinzip nahe.
\end{remark}

\begin{remark}[Komplexitätstheoretische Analogie]
\label{rem:pnp-analogy}
Das arithmetische Positivitäts-Framework besitzt ein komplexitätstheoretisches Analogon (vgl.\ das Begleitpaper zu P vs NP): Algebraische (Hodge-)Klassen entsprechen ,,uniformen Algorithmen'' (polynomialzeit-berechenbaren Zeugen), während transzendente Klassen ,,schweren Zeugen'' mit hoher Kolmogorov-Komplexität entsprechen. Das Versagen arithmetischer Positivität für eine transzendente Klasse ist analog zur Entropie-Separationsvermutung: Objekte, die sich strukturierter Beschreibung (algebraisch/algorithmisch) widersetzen, weisen eine Positivitäts-/Entropie-Lücke auf.
\end{remark}


\subsection{Richtungen für stärkere Positivität}
\label{sec:stronger_positivity}

Da $\mathrm{AP}^p = \mathrm{AbsHodge}^p$ (Bemerkung~\ref{rem:ap_equals_absolute}), gehen die aktuellen AP-Bedingungen nicht über Delignes Absolutheit hinaus. Dieser Abschnitt untersucht Kandidaten zur \textit{Verschärfung}, die genuin neue Einschränkungen jenseits der Hodge-Riemann-Relationen liefern könnten.

\subsubsection{\texorpdfstring{$p$-adische Steigungs-Constraints}{p-adische Steigungs-Constraints}}

\begin{remark}[$p$-adische Positivität: Newton- vs.\ Hodge-Polygon]
\label{rem:newton-hodge}
Sei $X/\mathbb{Q}$ eine glatte projektive Varietät mit guter Reduktion bei einer Primzahl $p$. Die kristalline Kohomologie $H^{2p}_{\mathrm{cris}}(X_p/\mathbb{Z}_p)$ trägt einen Frobenius-linearen Endomorphismus $\varphi$, dessen Eigenwert-Steigungen durch das Newton-Polygon beschränkt werden. Die Katz-Vermutung (bewiesen von Mazur~\cite{Mazur1972} und Ogus) besagt, dass das Newton-Polygon auf oder über dem Hodge-Polygon liegt.

Man definiere eine \textit{$v$-adische Steigungs-Constraint} (AP$'$): Für eine rationale Klasse der Kodimension $r$ soll an jeder guten endlichen Stelle $v$ die entsprechende kristalline beziehungsweise $v$-adische Realisierung im ,,algebraischen Steigungsbereich'' liegen, also Frobenius-Steigung $r$ besitzen. Algebraische Klassen erfüllen dies automatisch. Für nicht-algebraische absolute Hodge-Klassen könnte die Steigung von $r$ abweichen, was eine neue Obstruktion (O4) jenseits der Absolutheit liefern würde. Diese Formulierung trennt die Kodimension $r$ bewusst von der endlichen Stelle $v$.
\end{remark}

\subsubsection{Arakelov-Positivität}

\begin{remark}[Arakelov-Positivität als globaler Verstärker]
\label{rem:arakelov-positivity}
Die Arakelov-Schnittpaarung auf arithmetischen Varietäten kombiniert: (i) Schnittmultiplizitäten in den Fasern eines regulären Modells $\mathcal{X} \to \mathrm{Spec}(\mathbb{Z})$ (endliche Primstellen); (ii) den Green-Strom an der archimedischen Stelle, verwandt mit $Q_L$ über die Hodge-Metrik. Eine \textit{Arakelov-verstärkte AP-Bedingung} würde Positivität der \textit{globalen} Arakelov-Höhenpaarung $\hat{h}(\alpha, \alpha) \geq 0$ für alle Modelle erfordern, nicht nur der archimedischen Komponente. Da algebraische Zykel nicht-negativen Arakelov-Selbstschnitt haben (arithmetischer Hodge-Index-Satz von Moriwaki~\cite{Moriwaki2000}), ist dies ein natürlicher Kandidat für AP$'$.
\end{remark}

\begin{remark}[Testfall für AP$'$]
\label{rem:test-case-ap-prime}
Jede verstärkte Bedingung AP$'$ muss zwei Kriterien erfüllen: (1) \textit{Nicht-Trivialität:} AP$'$ darf \textit{nicht} automatisch von jeder absoluten Hodge-Klasse erfüllt werden. (2) \textit{Algebraische Sättigung:} Jede algebraische Klasse muss AP$'$ erfüllen. Ein Testfall ist eine Varietät $X$ mit einer absoluten Hodge-Klasse $\alpha$, deren Algebraizität vermutet aber nicht bewiesen ist. Kandidat: die Hodge-Klasse auf einer allgemeinen abelschen Vierfalte ($g = 4$, $p = 2$), die außerhalb des Lefschetz-Rings liegt, aber nach Delignes Satz absolut ist.
\end{remark}

\subsubsection{Tannakische Perspektive}

\begin{remark}[Tannakische Rekonstruktion für CM-Fälle]
\label{rem:tannakian}
Für eine abelsche Varietät $A$ vom CM-Typ ist die Mumford-Tate-Gruppe $\mathrm{MT}(A)$ ein Torus, und die Kategorie der von $H^1(A)$ erzeugten Hodge-Strukturen ist eine neutrale Tannakische Kategorie über $\mathbb{Q}$. Die Hodge-Klassen sind genau die $\mathrm{MT}(A)$-Invarianten. Die AP-Bedingungen lassen sich in der Sprache der Mumford-Tate-Gruppe umformulieren (vgl.\ Deligne--Milne~\cite{DeligneMilne1982}). Um über dies hinauszugehen, bräuchte man \textit{motivische} Positivitäts-Constraints, die nicht von der Mumford-Tate-Gruppe allein erfasst werden.
\end{remark}

\begin{remark}[Mumford-Tate-Analyse für Vermutung~\ref{conj:ghr}]
\label{rem:mumford-tate-ghr}
Für abelsche Varietäten $A$ liefert Delignes Satz
$H^{2p}(A,\mathbb{Q})\cap H^{p,p}(A)=\mathrm{AbsHodge}^p(A)=\mathrm{AP}^p(A)$.
Die GHR-Obstruktion kann eine hypothetische nichtalgebraische Hodge-Klasse auf einer abelschen Varietät daher nicht unterscheiden: falls eine solche Klasse existiert, ist sie bereits absolut und ihr GHR-Spektrum ist durch die Hodge-Riemann-Relationen nichtnegativ. Als Beweiswerkzeug ist die Obstruktion auf abelschen Varietäten also vakuos; die Hodge-Vermutung für abelsche Varietäten ist genau die verbleibende Aussage, dass diese absoluten/AP-Klassen algebraisch sind. Fortschritt hängt von der Mumford-Tate-Vermutung ab: Falls $\mathrm{MT}(A)$ treu auf $H^*(A, \mathbb{Q})$ wirkt und die $\mathrm{MT}(A)$-Invarianten algebraisch sind, folgt die Hodge-Vermutung für $A$. Dies ist bekannt für $\dim A \leq 5$ (Moonen~\cite{Moonen1999}) und für CM-Typ (Abdulali~\cite{Abdulali1994}).
\end{remark}

\subsubsection{Prismatische und kategoriale Ansätze}

\begin{remark}[Prismatische Positivität (AP4)]
\label{rem:prismatic}
Bhatts und Scholzes prismatische Kohomologie~\cite{BhattScholze2022} -- aufbauend auf der integralen $p$-adischen Hodge-Theorie von Bhatt--Morrow--Scholze~\cite{BhattMorrowScholze2018} -- liefert einen vereinheitlichten Rahmen für $p$-adische Kohomologietheorien. Eine hypothetische \textit{prismatische Positivitätsbedingung} (AP4) würde erfordern, dass die prismatische Realisierung $\alpha_{\Delta}$ in einem ,,positiven Kegel'' liegt, definiert durch Nygaard-Filtration und Frobenius-Aktion. Algebraische Klassen erfüllen dies automatisch. Nicht-algebraische absolute Hodge-Klassen könnten AP4 verletzen, was eine neue Obstruktion jenseits (O1)--(O3) liefert.
\end{remark}

\begin{remark}[Prismatische Kohomologie-Route]\label{rem:prismatic-route}
Nie \cite{Nie2021} konstruierte absolute Hodge-Zykel in der prismatischen Kohomologie abelscher Varietäten über $p$-adischen Körpern unter Verwendung von Andr\'es motivierten Zykeln im Bhatt--Scholze-Framework. Prismatische Kohomologie $R\Gamma_\Delta(X/S)$ interpoliert \'etale, de Rham- und kristalline Kohomologie integral und liefert einen universellen kohomologischen Rahmen. Die Verbindung zur Positivität über Bridgeland-Stabilitätsbedingungen auf derivierten Kategorien legt nahe, dass nichtlineare, derivierte Positivitätsbedingungen -- die die rigide Galois-Struktur prismatischer Kristalle mit geometrischer Stabilität kombinieren -- die Lücke zwischen Absolutheit und Algebraizität überbrücken könnten.
\end{remark}

\subsubsection*{Schwellenaxiom: jenseits von Deligne}

\begin{axiom}[Sprachbrücke -- JD-Hodge]\label{ax:jd-hodge}
Es existiert eine Kategorie $\mathcal{C}$ von ,,Hodge-Trägern'' (Objekte, die kohomologische Daten repräsentieren), ausgestattet mit:
\begin{enumerate}
\item[(JD1)] \textbf{Funktorialität:} Rückzug, Vorwärtsbild und Produkte wirken auf $\mathcal{C}$ ohne nicht-kanonische Wahlen.
\item[(JD2)] \textbf{Intrinsische Filtration:} Eine funktorielle Hodge-Filtration $F^\bullet$ ist kanonisch definiert (nicht über eine Wahl von Einbettung oder Modell).
\item[(JD3)] \textbf{Kanonische Polarisation:} Eine mit (JD1) und (JD2) kompatible Positivitätsstruktur existiert, die die Hodge--Riemann-Form verallgemeinert.
\end{enumerate}
In dieser Sprache würde die Hodge-Vermutung nur dann zu einer \emph{formalen Konsequenz}, wenn $\mathcal{C}$ zusätzlich ein Zyklus-Realisierungstheorem trägt: Positivität in $\mathcal{C}$ müsste Zugehörigkeit zum Bild der klassischen Zyklusklassenabbildung implizieren. Delignes Absolutheit wäre dann ein notwendiger Schatten der reicheren Struktur in $\mathcal{C}$, nicht selbst der fehlende Algebraizitätsschritt.
\end{axiom}

\begin{remark}[Der Nicht-Kanonizitäts-Scanner]\label{rem:noncanonicity-scanner}
Um zu lokalisieren, wo die aktuelle Theorie hinter JD-Hodge zurückbleibt, suche nach drei Signaturen:
\begin{enumerate}
\item \textbf{Wahlabhängigkeit:} Wo treten Spaltungen oder Vergleichsisomorphismen (de Rham/Betti/$\ell$-adisch) auf, die ,,existieren'', aber keine natürlichen Transformationen sind?
\item \textbf{Inkompatible Operationen:} Welche geometrische Operation (Rückzug unter nicht-endlichen Morphismen, Gysin-Vorwärtsbild) zerstört die Polarisation oder destabilisiert die Filtration?
\item \textbf{Vergleichsbrüche:} Wo hängt die Aussage von einem Vergleich zwischen Realisierungen ab, der nicht strukturell kontrolliert ist?
\end{enumerate}
Im aktuellen AP-Framework tritt der Bruch bei AP2a (Absolutheit) auf: Der Vergleichsisomorphismus $H^*_{\mathrm{dR}}(X) \otimes \mathbb{C} \cong H^*_{\mathrm{Betti}}(X^\sigma) \otimes \mathbb{C}$ ist \emph{kein} Morphismus in einer natürlichen Kategorie -- er erfordert die Wahl von $\sigma \in \mathrm{Aut}(\mathbb{C}/\mathbb{Q})$.
\end{remark}

\begin{remark}[Kandidaten für $\mathcal{C}$]\label{rem:candidates-C}
Drei Kandidaten für die Kategorie $\mathcal{C}$ werden derzeit untersucht:
\begin{enumerate}
\item \textbf{Prismatische Kohomologie} (Bhatt--Scholze, Nie \cite{Nie2021}): $R\Gamma_\Delta(X/S)$ interpoliert alle klassischen Realisierungen integral, mit einer kanonischen Frobenius-Aktion für (JD2). Die Positivität (JD3) ist noch nicht formuliert, aber strukturell möglich via Bridgeland-Stabilität.
\item \textbf{Motivische Kohomologie} (Voevodsky): Die derivierte Kategorie der Motive $\mathrm{DM}(k)$ erfüllt (JD1) per Konstruktion. Die Standard-Vermutungen würden (JD2) und (JD3) liefern, sind aber unbewiesen.
\item \textbf{Gemischte Hodge-Moduln} (Saito): Erfüllen (JD1)--(JD2) für glatte projektive Varietäten, aber (JD3) ist auf die klassische HR-Form beschränkt -- genau die AP$=$AbsHodge-Barriere dieses Papers.
\end{enumerate}
\end{remark}

\begin{remark}[Zyklus-Realisierungs-Gate für abelsche Varietäten]\label{rem:jd-abelian}
Für abelsche Varietäten macht Delignes Satz jede Hodge-Klasse absolut und damit AP nach Proposition~\ref{prop:absolute_implies_ap}. Eine JD-Hodge-Kategorie würde dies nur verbessern, wenn sie ein zusätzliches Zyklus-Realisierungs-Gate liefert: AP-/absolute Klassen, die zu kanonischen Morphismen in $\mathcal{C}$ werden, müssten nachweislich im Bild der klassischen Zyklusklassenabbildung liegen. Ohne dieses Zusatztheorem ersetzt man ,,absolut'' nur durch ,,motivisch'' oder ,,kanonisch in $\mathcal{C}$'' und formuliert das verbleibende Algebraizitätsproblem neu. Der abelsche Fall ist daher ein Stresstest für JD-Hodge, keine bereits aus (JD1)--(JD3) gewonnene Konsequenz.
\end{remark}

\begin{remark}[Die Deligne-Grenze als Sprachdefekt]\label{rem:deligne-language}
Axiom~\ref{ax:jd-hodge} identifiziert die Obstruktion als \emph{linguistisch}, nicht rechnerisch: Die Vergleichsisomorphismen $\mathrm{comp}_{\mathrm{dR} \leftrightarrow B}$ und $\mathrm{comp}_{\ell\text{-adic} \leftrightarrow B}$ existieren, sind aber nicht funktoriell in $X$. Absolute Hodge-Klassen sind genau diejenigen, die unter allen Realisierungen kompatibel sind -- und Algebraizität wird zu $\mathrm{AH}(X) = \mathrm{Im}(\mathrm{CH}^*(X) \to H^*(X))$. Wären die Vergleiche natürliche Transformationen in einer Kategorie $\mathcal{C}_{\mathrm{abs}}$, wäre dies eine formale Konsequenz. Die Deligne-Grenze lautet daher: \emph{Nicht-kanonische Spaltungen machen jeden Algebraizitätsbeweis instabil unter Funktoren}.
\end{remark}

\begin{remark}[Bridgeland-Stabilität und Positivität auf derivierten Kategorien]
\label{rem:bridgeland}
Bridgelands Stabilitätsbedingungen~\cite{Bridgeland2007} auf der derivierten Kategorie $D^b(X)$ liefern einen Positivitätsbegriff für Objekte statt für Kohomologieklassen. Eine mögliche Verschärfung: Fordere, dass die Kohomologieklasse $\alpha$ als Chern-Charakter eines Bridgeland-stabilen Objekts entsteht. Dies ist strikt stärker als eine Hodge-Klasse zu sein.
\end{remark}

\begin{remark}[Spektrale Lücke auf Vektorbündeln]
\label{rem:spectral-gap}
Man betrachte die ,,Spektraldaten'' eines geeigneten Differentialoperators (z.B.\ des Laplace-Operators auf $\bigwedge^p \Omega_X$). Algebraische Zykel entsprechen Nullmoden (harmonischen Formen), während nicht-algebraische Hodge-Klassen ,,fast-Nullmoden'' mit einer kleinen, aber nicht-verschwindenden Spektrallücke entsprechen würden. Eine Positivitätsbedingung basierend auf dieser Spektrallücke könnte algebraische von nicht-algebraischen Klassen unterscheiden.
\end{remark}

\subsubsection{Geometrische und o-minimale Ansätze}

\begin{remark}[VHS-Familien und Monodromie]
\label{rem:vhs-monodromy}
Für eine Variation von Hodge-Strukturen (VHS) über einer Basis $S$ zeigt der Satz von Cattani--Deligne--Kaplan~\cite{CattaniDeligneKaplan1995}, dass Hodge-Loki abzählbare Vereinigungen geschlossener algebraischer Untervarietäten von $S$ sind. Das ist eine Aussage darüber, wo flache rationale Klassen vom Hodge-Typ bleiben, nicht darüber, dass die entsprechenden Klassen algebraisch sind. Eine verstärkte AP-Bedingung könnte erfordern, dass $\alpha_s$ sich als \textit{flacher Schnitt} über eine Zariski-dichte Teilmenge von $S$ fortsetzt und zusätzlich unabhängige Zyklus-Provenienz besitzt. Dies hängt mit der Halbeinfachheit der Monodromie-Darstellung und mit der von der Hodge-Vermutung erwarteten Algebraizität flacher Schnitte zusammen.
\end{remark}

\begin{remark}[$p$-adische Positivität via O-Minimalität]
\label{rem:o-minimality}
Die Theorie o-minimaler Strukturen (Bakker--Brunebarbe--Tsimerman~\cite{BBT2023}) wurde angewendet, um Definierbarkeitsresultate für Periodenabbildungen und Hodge-Loki zu beweisen. Eine mögliche AP$'$-Bedingung: Fordere, dass die ,,Umlaufbahn'' von $\alpha$ unter der definierbaren Periodenabbildung algebraisch ist. Nach CDK gilt dies für algebraische Klassen. Für nicht-algebraische Hodge-Klassen könnte die Umlaufbahn ,,wild'' (transzendent) sein, was die Zahmheitsbedingung verletzt. Dies steht in Verbindung mit der Arbeit von Binyamini--Novikov~\cite{BinyaminiNovikov2017} zum Zählen rationaler Punkte auf definierbaren Mengen.
\end{remark}

\begin{remark}[Galois-Defekte für transzendente Klassen]
\label{rem:galois-defects}
Wenn eine Hodge-Klasse $\alpha$ nicht absolut ist, existiert $\sigma \in \mathrm{Aut}(\mathbb{C})$ mit $\alpha^\sigma \notin H^{p,p}(X^\sigma)$. Die ,,Größe'' dieses Defekts -- gemessen etwa an der Komponente von $\alpha^\sigma$ in $H^{p-1,p+1}(X^\sigma) \oplus H^{p+1,p-1}(X^\sigma)$ -- ist ein quantitatives Maß für Nicht-Algebraizität. Wächst dieser Defekt unter Iteration der Galois-Konjugation, würde das GHR-Spektrum negativ werden, was eine berechenbare Obstruktion liefert.
\end{remark}

\begin{remark}[Ricci-artiger Fluss auf $(p,p)$-Formen]
\label{rem:ricci-flow}
In Analogie zum Ricci-Fluss betrachte man einen Fluss auf $(p,p)$-Formen:
$\partial \alpha / \partial t = -\Delta_L \alpha + \lambda \alpha$,
wobei $\Delta_L$ der Lefschetz-Laplace-Operator ist. Algebraische Klassen -- als harmonische Repräsentanten algebraischer Zykel -- sind Fixpunkte (,,thermodynamische Minima'') dieses Flusses. Nicht-algebraische Hodge-Klassen wären unter dem Fluss \textit{instabil}. Dies interpretiert algebraische Zykel als Attraktoren einer natürlichen geometrischen Evolution.
\end{remark}

\subsection{Jenseits bilinearer Positivität: nichtlineare und gemittelte Bedingungen}
\label{sec:nonlinear-positivity}

Die Hodge--Riemann-Relationen liefern \emph{bilineare} Positivität: $(-1)^p Q_L(\alpha, \bar\alpha) > 0$ für primitive $(p,q)$-Klassen. Wie in den Einschränkungen (Abschnitt~9) angemerkt, genügt dies nicht für die volle Hodge-Vermutung. Wir identifizieren drei Richtungen für \emph{nichtlineare} oder \emph{gemittelte} Positivität, die über HR hinausgehen:

\begin{remark}[Drei Wege zu nichtlinearer Positivität]
\label{rem:nonlinear-positivity}
\emph{Rand-Positivität.}
Für eine Familie ampler Klassen $L_t \to L_0$, wobei $L_0$ nef aber nicht ample ist, können die vorzeichenkorrigierten Werte $(-1)^p Q_{L_t}$ für $t \to 0$ nichtnegative Grenzwerte haben, selbst wenn $Q_{L_0}$ nicht definiert ist. Diese ,,Rand-Positivität'' erfasst Klassen, die im Grenzwert algebraisch sind, aber von keiner einzelnen Polarisierung detektiert werden.

\emph{Positivität unter motivischer Faltung.}
Für eine Produktvarietät $X \times Y$ liefert die Künneth-Zerlegung diagonale Summanden $H^{a,a}(X)\otimes H^{p-a,p-a}(Y)$, auf denen die Vorzeichenrechnung für Produktpolarisierungen multiplikativ ist (Lemma~\ref{lem:sign_product}). Diese Aussage ist jedoch nur ein kontrollierter Produktpolarisierungs-Test. Gemischte Hodge-Typen und primitive Kreuzterme auf $X\times Y$ werden dadurch nicht beherrscht; ihre Positivitätsanalyse müsste über den diagonalen Rahmen dieser Arbeit hinausgehen.

\emph{Positivität nach Galois-Mittelung.}
Für eine absolute Hodge-Klasse $\alpha$ ist ein vorsichtig definierter endlicher Galois-Mittelwert $\bar\alpha$ (vgl.\ Bemerkung~\ref{rem:galois-profinite}) wieder eine absolute Hodge-Klasse. Für $p=1$ sind solche gemittelten Klassen nach dem Lefschetz-$(1,1)$-Theorem algebraisch. Für $p\geq2$ ist die Algebraizität des Mittelwerts im Allgemeinen \emph{nicht bekannt}; sie ist eine Spezialform von Delignes Frage. Die arithmetischen Positivitätsbedingungen AP1--AP3 lassen sich allenfalls als Nähe-Heuristik zu diesem Mittelwert in der $Q_L$-Metrik lesen. Diese Umformulierung ist nur dann nützlich, wenn der Mittelwert unabhängig von der Hodge-Vermutung als algebraisch gezeigt werden kann.
\end{remark}

\begin{remark}[Galois-Mittelung: profiniter Limes]\label{rem:galois-profinite}
Die naive Mittelung $\bar{\alpha} = \frac{1}{|\mathrm{Aut}(\mathbb{C}/\mathbb{Q})|} \sum_\sigma \sigma^*\alpha$ ist nicht wohldefiniert, da $\mathrm{Aut}(\mathbb{C}/\mathbb{Q})$ eine profinite Gruppe der Kardinalität $2^{\mathfrak{c}}$ ist. Eine vorsichtige Formulierung arbeitet nur auf endlicher Stufe: Für jede endliche Galois-Erweiterung $K/\mathbb{Q}$, über der Modell und Vergleichsdaten kontrolliert sind, bilde man
\[
\bar{\alpha}_K = \frac{1}{|\mathrm{Gal}(K/\mathbb{Q})|}
  \sum_{\sigma \in \mathrm{Gal}(K/\mathbb{Q})} \sigma^*\alpha.
\]
Jede projektive-Limes-Deutung braucht zusätzliche Abstiegs- und Kompatibilitätsdaten; sie darf nicht automatisch als algebraische Klasse behandelt werden.
\end{remark}

\begin{remark}[Universeller Selektionsmechanismus]
\label{rem:universal-selection}
Alle drei Wege teilen die in der Meta-Analyse des FST-Programms identifizierte Struktur: \emph{lokale Positivität + Regularisierung/Mittelung $\Rightarrow$ globale Selektion}. Die ,,Regularisierung'' nimmt verschiedene Formen an:
\begin{itemize}
\item NS: Zeitmittelung der Nächstpunkt-Selektion,
\item DE: RG-Fluss vom UV- zum IR-Fixpunkt,
\item BSD: Scheibenmittelung von $L$-Funktionswerten,
\item Hodge: Galois-Mittelung oder Polarisierungs-Degeneration.
\end{itemize}
Dieser strukturelle Isomorphismus legt nahe, dass ein einziges abstraktes Theorem -- ein ,,Selektion durch Positivität''-Prinzip -- allen Instanziierungen zugrunde liegt.
\end{remark}

\subsection{Numerische Illustration}

Das Python-Skript \texttt{compute\_ghr\_spectrum.py} (dieser Arbeit beigelegt) verifiziert die AP-Bedingungen numerisch an 17 Testfällen:
\begin{itemize}
\item K3-Flächen (4 Polarisierungen), abelsche Flächen (5 Polarisierungen), abelsche Vierfache ($p = 2$), CM-Endomorphismenklassen auf $E \times E$;
\item Voisin-artige Beispiele: Fermat-Quintik, $K3 \times K3$-Tensorprodukte, Weil-Torus;
\item Systematischer Polarisierungs-Scan: 500 zufällige ample Klassen auf einem $(1, 19)$-Signatur-Modell;
\item Simulierte Galois-Aktionen: 50 orthogonale Rotationen des negativen Eigenraums;
\item SDP-Kegelstruktur: Eigenwert-Landschaft von $M(L) = (-1)^p Q_L|_{P^{p,p}}$ über dem Amplekegel.
\end{itemize}

Da $\mathrm{AP}^p = \mathrm{AbsHodge}^p$ ein Theorem ist und keine empirische Vermutung, dienen die numerischen Rechnungen nur als pädagogische Illustration. Alle 17 Testkonfigurationen zeigen $(-1)^p Q_L \geq 0$ auf $P^{p,p}$, wie es die Hodge-Riemann-Relationen vorgeben.

\begin{remark}[Voisin-GHR-Obstruktionstest]
\label{rem:voisin-ghr}
Die GHR-Obstruktion wurde an sechs Varietäten illustriert: generischer $T^4$, Weil-Torus (CM durch $\mathbb{Z}[i]$), $K3 \times K3$, Fermat-Quintik $X_5$, Voisin-Typ nicht-projektiver Torus und eine SDP-Relaxierung des AP-Kegels. Für alle projektiven Beispiele ist AP1 durch die Hodge--Riemann-Relationen garantiert. Der $K3 \times K3$-Test mit Nicht-Tensorprodukt-Störungen erzeugt 3 negative Eigenwerte von $(-1)^p Q_L$; dies illustriert Detektion von Nichtabsolutheit bzw. von Modellklassen, die nach Konjugation nicht im zulässigen $(p,p)$-/Primitivrahmen bleiben. Es ist kein Algebraizitätskriterium. Die AP-Kegel-Analyse zeigt alle zufälligen Klassen innerhalb des Kegels mit positiver Marge ($0{,}24$), was nur die Stabilität der Illustration anzeigt. Skript: \texttt{compute\_voisin\_test.py}.
\end{remark}

\subsection{Kontrollierte Testfamilien jenseits von \texorpdfstring{$\mathrm{AP}=\mathrm{AbsHodge}$}{AP=AbsHodge}}
\label{subsec:zookeeper-hodge-transfer}

Die Post-Zookeeper-Deutung dieses Papers ist bewusst bescheiden.  Da
arithmetische Positivität auf absolute Hodge-Klassen kollabiert, ist der
nächste Schritt kein weiteres globales Positivitätsaxiom, sondern eine
Testfamilien-Leiter für mögliche Verstärkungen jenseits von Delignes
Absolutheit.
Die Symbole $\mathrm{AP}'$, $\mathrm{AP}''$ und
$\mathrm{AP}^{\prime\prime\prime}$ bezeichnen keine globalen
Ersatzvermutungen für die Hodge-Vermutung. Sie bezeichnen zunehmend
strengere Tests auf kontrollierten Spezialfamilien: $\mathrm{AP}'$ erfasst
Provenienzdefekte, $\mathrm{AP}''$ fordert ein beschränktes
Orbit-Rigiditätszertifikat, und $\mathrm{AP}^{\prime\prime\prime}$ filtert
Zertifikate aus, deren Zusatzwahlen lediglich die Zielklasse kodieren. Keine
dieser Bedingungen wird als allgemeiner Beweis der Hodge-Vermutung behauptet.

Die Notation ist daher eine Claim-Status-Konvention, kein Satzname. Jede
Algebraizitätsaussage in diesem Abschnitt ist nur so stark wie die benannten
Familienannahmen, die unabhängig von der Zielklasse $\alpha$ verifiziert
werden. Werden die Zusatzdaten erst nach Kenntnis von $\alpha$ gewählt, gilt
die Bedingung als gescheiterter Test und nicht als Evidenz.

Die aktuelle CORE-Schließung von Riemann-$C2/MS2$ verstärkt das
Hodge-Resultat dieses Papers nicht.  Hodge bleibt ein negativer
Kontrollfall des Programms: Hodge--Riemann-Positivität sättigt zu breit,
sodass jeder künftige Transfer wirklich neue arithmetische oder motivische
Zusatzbedingungen braucht und nicht die Riemann-Kernel-Schließung
wiederverwenden darf.

Die fünf Transferfelder lauten:
\begin{description}
  \item[Operator/Form.] Hodge-Laplacian, Mumford--Tate- und
  Monodromieoperatoren sowie die Periodenabbildung.
  \item[Defekt.] Der transzendente Restanteil einer absoluten Hodge-Klasse
  nach Abzug der in der Familie bekannten algebraischen Zykelspanne.
  \item[Approximation.] Abelsche Vierfalten vom Weil-Typ, singuläre
  OG6-hyperkählersche Modelle, Potenzen bekannter Beispiele und
  Reduktionen in Charakteristik $p$.
  \item[Gap.] Zusatzbedingungen jenseits der Hodge--Riemann-Positivität:
  $p$-adische Slopes, Tate-Constraints, Mumford--Tate-Restriktionen und
  o-minimale Algebraizität.
  \item[Grenzübergang.] Deformation, Potenzen und Spezialisierung von
  kontrollierten Familien in einen breiteren motivischen Rahmen.
\end{description}

Eine nützliche nächste Version sollte deshalb eine Tabelle von
Testfamilien enthalten.  Für jede Familie werden festgehalten: bekannter
Hodge-/Tate-Status, der in der Literatur verwendete Zykelmechanismus, die
vorgeschlagene stärkere Bedingung $\mathrm{AP}'$ und ob der Defekt
verschwindet.  Weil-Vierfalten mit Diskriminante eins, OG6-Konstruktionen,
abelsche Varietäten über endlichen Körpern mit Standardvermutung vom
Hodge-Typ und Voisin-artige Stresstests bilden die natürliche erste
Leiter.  Wenn $\mathrm{AP}'$ auch in allen Negativkontrollen automatisch
gilt, ist sie zu schwach; wenn sie eine bekannte positive Familie von
einem Stresstest trennt, wird sie zu einem echten Kandidaten für weitere
Formalisierung.

\subsubsection*{Hin zu \texorpdfstring{$\mathrm{AP}'$}{AP'}: Provenienzdefekte}
\label{subsec:provenance-defects}

Die erste Verschärfung ist als Provenienzscanner zu lesen, nicht als ein
weiterer Positivitätskegel. Relativ zu einer benannten Testfamilie definiert
man einen Defektvektor
\[
D_{\mathrm{prov}}(\alpha)
  = (D_{\mathrm{slope}},D_{\mathrm{Tate}},D_{\mathrm{MT}},
     D_{\mathrm{mon}},D_{\mathrm{cycle}}),
\]
dessen Einträge Steigungs-Kompatibilität an guten Reduktionen, Erklärung
durch den Tate-/Galois-Fixteil, Mumford--Tate-Kompatibilität,
monodromische oder o-minimale Persistenz und die nach Abzug der bekannten
algebraischen Zykelspanne verbleibende Restklasse erfassen. Eine bedingte
Fingerprint-Aussage würde besagen, dass in einer Testfamilie mit unabhängig
verfügbaren Tate-/Standardvermutungsdaten und Vergleichsfunktorialität aus
$D_{\mathrm{prov}}(\alpha)=0$ folgt, dass $\alpha$ in der bekannten
algebraischen Zykelspanne liegt. Das ist eine Spezialfamilien-Diagnostik,
kein globales Theorem.

Jede Komponente von $D_{\mathrm{prov}}$ muss mit dem Status
\emph{verifiziert}, \emph{gescheitert} oder \emph{unbekannt} geführt
werden. Unbekannte Komponenten zählen nicht als verschwindende Defekte.
Diese Konvention verhindert, dass der Provenienzscanner tautologisch wird:
$\mathrm{AP}'$ gilt nur als bestanden, wenn alle erforderlichen Komponenten in
der gewählten Testfamilie unabhängig verifiziert sind.

\subsubsection*{Hin zu \texorpdfstring{$\mathrm{AP}''$}{AP''}: Bedingte Orbit-Rigidität}
\label{subsec:orbit-rigidity}

Die fünf Transferkanäle legen eine prototypische Verschärfung nahe.  Sei
$f\colon \mathcal{X}_{S}\to S$ eine Ausbreitung von $X$ über einer glatten
Basis~$S$, sei $Z\subset S$ eine irreduzible algebraische Komponente des
Hodge-Lokus von~$\alpha$, sei $B$ die zugehörige relative Chow- oder
Hilbert-Komponente, und sei $T\subset Z$ eine Zariski-dichte Teilmenge,
an deren Punkten $\alpha_{t}$ relativ zu~$B$ algebraisch ist.
Wir nennen $\alpha$ \emph{orbit-rigide auf $(f,Z,B)$}, wenn:
\begin{enumerate}
\item[(i)] \emph{Properness.} Die relativen Chow-/Hilbert-Komponenten,
  die die relevanten Zykelklassen parametrisieren, sind proper über~$S$.
\item[(ii)] \emph{Dichtheit.} $T$ ist Zariski-dicht in~$Z$.
\item[(iii)] \emph{Vergleichsfunktorialität.}  Die Vergleichsisomorphismen
  zwischen Betti-, de~Rham-, $\ell$-adischer und kristalliner
  Realisierung sind funktoriell auf dem relevanten Quotienten; beim Liften
  von einer speziellen zur generischen Faser entstehen keine neuen
  Restklassen.
\item[(iv)] \emph{Spezialisierungsrigidität.}  Basiswechsel entlang
  $Z\to S$ erzeugt keine neuen residualen Zykelklassen.
\end{enumerate}

\begin{lemma}[Bedingte Orbit-Rigiditäts-Schablone]
\label{lem:orbit-rigidity}
Sei $\alpha\in\mathrm{AbsHodge}^{p}(X)$ orbit-rigide auf $(f,Z,B)$
im obigen Sinne.  Dann liegt die Einschränkung $\alpha_{\eta}$ auf die
generische Faser $\mathcal{X}_{\eta}$ über dem generischen Punkt $\eta$
von~$Z$ im Bild der Zykelklassenabbildung
$\mathrm{cl}\colon\mathrm{CH}^{p}(\mathcal{X}_{\eta})_{\mathbb{Q}}\to
H^{2p}(\mathcal{X}_{\eta},\mathbb{Q}(p))$.
Liegt der Basispunkt $0\in S$ in derselben kontrollierten Komponente,
so ist $\alpha$ algebraisch.
Dies ist eine Checklisten-Aussage: Die Zykelprovenienz-, Vergleichs- und
Spezialisierungshypothesen gehören zum Input und werden nicht durch
$\mathrm{AP}$, $\mathrm{AP}'$ oder $\mathrm{AP}^{\prime\prime\prime}$ allein
verifiziert.
\end{lemma}

\begin{proof}[Beweisidee]
Properness von $B$ \emph{(i)} und Zariski-Dichtheit von $T$ \emph{(ii)}
implizieren über das Bewertungskriterium, dass der Abschluss des
algebraischen Lokus in $B$ den generischen Punkt~$\eta$ von~$Z$ enthält.
Vergleichsfunktorialität \emph{(iii)} stellt sicher, dass jede algebraische
Klasse über $t\in T$ konsistent in allen kohomologischen Realisierungen
liftet.  Spezialisierungsrigidität \emph{(iv)} verhindert die Akkumulation
transzendenter Restklassen am Limes.  Ein vollständiger Beweis erfordert
das Cattani--Deligne--Kaplan-Prinzip für Hodge-Loki
\cite{CattaniDeligneKaplan1995} zusammen mit einem Properness-Argument
für das relative Hilbert-Schema; beide Schritte sind bedingt durch die
Standardvermutungen in allen relevanten Charakteristik-$p$-Reduktionen oder
müssen durch unabhängige familienbezogene Zykelprovenienz-Eingaben ersetzt
werden.
\end{proof}

\begin{remark}
\label{rem:ap-double-prime}
Lemma~\ref{lem:orbit-rigidity} beweist nicht die Hodge-Vermutung.
Es ist als bedingtes Schließungsschema zu lesen, nicht als neues
unbedingtes Algebraizitätstheorem. Seine Aussagekraft liegt vollständig in
der unabhängigen Verifikation beschränkter relativer Chow-/Hilbert-Daten,
der Vergleichsfunktorialität und der Spezialisierungsrigidität in einer
benannten Testfamilie. Unter diesen separat geprüften Voraussetzungen
propagiert Algebraizität von einer Zariski-dichten Menge zum generischen
Punkt.  Weil-abelsche Vierfalten mit Diskriminante eins sind ein
natürlicher erster Kandidat für einen Positivtest; $K3\times K3$-Störungen
außerhalb des Tensorprodukt-Lokus dienen als erste Negativkontrolle.
Ob $\mathrm{AP}''$ diese beiden Familien trennt, entscheidet darüber,
ob es eine echte Verschärfung jenseits von $\mathrm{AbsHodge}$ darstellt.
\end{remark}

\subsubsection*{Hin zu \texorpdfstring{$\mathrm{AP}^{\prime\prime\prime}$}{AP'''}: Uniformität und Anti-Advice}
\label{subsec:ap-triple-prime}

Das Orbit-Rigiditätszertifikat besitzt weiterhin einen Fehlermodus: Die
Familie $f$, die Hodge-Lokus-Komponente $Z$, die Chow-/Hilbert-Komponente
$B$ oder die dichte Menge $T$ könnten erst nach Kenntnis der Zielklasse
$\alpha$ gewählt werden. Dann kann das Zertifikat $\alpha$ kodieren, statt
es zu erklären. Wir verwenden daher
\[
\mathrm{AP}^{\prime\prime\prime}_{\mathrm{unif}}(\alpha;R,C)
\]
als Anti-Tautologie-Filter für ein $\mathrm{AP}''$-Zertifikat $C$. Gefordert
ist, dass die Zusatzdaten aus einer vorregistrierten Regelklasse $R$ erzeugt
werden, die nur von den ambienten Daten $(X,p)$ oder von einer festen
Testfamilienleiter abhängt, mit beschränkter Zusatzinformation, Erosionsstabilität
nach Entfernen einer $\alpha$-spezifischen Inzidenzbedingung, Träger auf
einer Nachbarfamilie oder Zariski-dichten Menge und Scheitern auf
Negativkontrollen, die nur durch Hodge--Riemann-Formstabilität getrieben
sind. Diese Bedingung fügt keine neue globale Hodge-Vermutung hinzu; sie
sagt, wann ein Spezialfamilien-Zertifikat nicht tautologisch ist.

Operativ scheitert ein Zertifikat am
$\mathrm{AP}^{\prime\prime\prime}$-Filter, sobald einer der folgenden Tests
scheitert:
\begin{enumerate}
\item \emph{Vorregistrierung:} Die Regelklasse $R$ wurde nicht vor der
  Zielklasse fixiert oder folgt nicht aus ambienten Daten.
\item \emph{Beschränkte Zusatzinformation:} Die Wahlen von $f,Z,B,T$ tragen genug
  Information, um $\alpha$ zu rekonstruieren.
\item \emph{Erosionsstabilität:} Das Entfernen einer $\alpha$-spezifischen
  Inzidenzbedingung zerstört das Zertifikat.
\item \emph{Nachbarschaftsschwarm:} Dieselbe Regel trägt keine
  Nachbarfamilie und keine Zariski-dichte Menge.
\item \emph{Negativkontrollen:} Voisin-/$K3\times K3$-Stressfamilien
  bestehen nur wegen Hodge-Riemann-Formstabilität.
\end{enumerate}

\begin{center}
\resizebox{\columnwidth}{!}{%
\begin{tabular}{lllllll}
\toprule
\textbf{Familie} & \textbf{Zykelmechanismus} & \textbf{AP$'$} & \textbf{AP$''$} & \textbf{AP$^{\prime\prime\prime}$} & \textbf{Status} & \textbf{Ledger} \\
\midrule
CM-/Weil-Vierfalten & Endomorphismuszyklen & Slope/Tate/MT & Kandidat & Regel aus Endomorphismen & geplanter Positivkontroll-Kandidat & ungefüllt \\
Abelsche Vierfalten & bekannte Zykelspanne & Tate/MT & Kandidat & ambiente Familienregel & geplanter Positivkontroll-Kandidat & ungefüllt \\
OG6-Typ-Familien & hyperkählersche Zykel & Monodromie/Zykel & offen & Registry nötig & Testfall & geplant \\
Char-$p$-Tate-Fälle & Tate-Klassen & Slope/Tate & bedingt & Reduktionsregel & arithmetischer Test & geplant \\
Voisin/$K3\times K3$-Stress & nicht fixiert & Restklasse & erwartetes Scheitern & Post-hoc-Risiko & geplanter Stresskontroll-Kandidat & ungefüllt \\
\bottomrule
\end{tabular}}
\end{center}

\subsubsection*{Skalen- und Multi-Erosions-Guardrails}
\label{subsec:scale-reef}

Der $\mathrm{AP}^{\prime\prime\prime}$-Filter braucht zwei praktische Guardrails, bevor er als ernstes Testfamilien-Protokoll nutzbar ist.

\emph{Skalenspezifität.} Ein Provenienzdetektor sollte ein Skalen-Ledger führen:
\[
\begin{gathered}
\mathrm{scale\_level},\quad \mathrm{detector\_degree},\\
\mathrm{target\_complexity\_degree}
\end{gathered}
\]
für die untersuchte Familie. Ein Detektor, der nur eine niedergradige Formsignatur sieht, ist keine Provenienzbrücke, wenn die auszuschließende residuale nichtalgebraische Stressfamilie höhere Komplexität besitzt. Dann lautet der korrekte Status Spezifitäts- oder Kapazitätslücke, nicht Fortschritt in Richtung Algebraizität. Operativ sollte die Provenienzmatrix einen Restterm
\[
D_{\mathrm{scale}}(\alpha)
\]
ergänzen, der nur dann verschwindet, wenn der Detektor dieselbe Komplexitätsstufe auflöst wie die Klasse, die er ausschließen soll.
Operativ bedeutet das: $D_{\mathrm{scale}}(\alpha)=0$ ist nur zulässig, wenn der dokumentierte Detektorgrad mindestens dem dokumentierten Zielkomplexitätsgrad entspricht. Andernfalls lautet der Status Spezifitätslücke. Insbesondere sind Hodge-Riemann-Formsignaturen niedergradige Diagnostik; allein erklären sie K3$\times$K3- oder Voisin-artige Stressfamilien nicht als Provenienzdaten.

\emph{Multi-Erosions-Provenienz.} Für eine vorregistrierte Menge unabhängiger Erosionsachsen
\[
E=\{e_1,\ldots,e_m\},
\]
schreiben wir
\[
\mathrm{Reef}_{E}(\alpha;R,C,E)
\]
für folgende bedingte Erweiterung des Anti-Advice-Filters. Die Achsen $e_i$ werden aus ambienten Familiendaten fixiert, nicht aus den Koordinaten von $\alpha$; jedes Rand- oder Residuenbild $\mathrm{Res}_{e_i}(\alpha)$ wird entweder durch eine feste algebraische Zykelspanne auf dem jeweiligen Stratum erklärt oder als Defekt $D_{\mathrm{res}}(e_i,\alpha)$ geführt; Residuen stimmen auf Überlappungen im Cech-/Gysin-/nearby-cycle-Sinn überein; und die globale Rückhub-Obstruktion
\[
\mathrm{Ob}_{\mathrm{lift}}(\alpha;E)=0
\]
verschwindet. Das ist kein neuer Satz. Es ist eine Folgearbeits-Checkliste: Ein Spezialfamilien-Zertifikat soll scheitern, wenn dieselbe Klasse durch private Erosionsachsen, inkompatible Rand-Erklärungen oder eine ungelöste Abel-Jacobi-/Normalfunktions-/vanishing-cycle-Rückhub-Obstruktion nachgebaut werden kann.
Mindestens zwei unabhängige vorregistrierte Achsen sollen verwendet werden; eine einzelne maßgeschneiderte Achse ist nur eine Sonde, kein Riff-Zertifikat. Unbekannter Residuen-, Overlap-, Rückhub- oder Negativkontrollstatus ist ein offener Defekt und kein bestandener Test.

\subsection*{Ergebnisklassifikation}

Diese Arbeit ist ein \textbf{Grenzresultat}: Sie bestimmt präzise, was positivitätsbasierte Methoden für die Hodge-Vermutung leisten können und was nicht.

\begin{center}
\resizebox{\columnwidth}{!}{%
\begin{tabular}{lll}
\toprule
\textbf{Ergebnis} & \textbf{Typ} & \textbf{Verwendung im Paper} \\
\midrule
\multicolumn{3}{l}{\textit{No-Go-/Grenzresultat:}} \\
AP $=$ AbsHodge (HR liefert nichts Neues) & Satz & bewiesene Grenze \\
APV-Hard-Direction ist Delignes Frage & Korollar & Reichweitenbegrenzung \\
\midrule
\multicolumn{3}{l}{\textit{Konstruktive Werkzeuge:}} \\
GHR-Spektrum-Definition & Definition & nur Visualisierung \\
Endlicher Rückzug / Künneth-Vorzeichentest & Partiell, bedingt & begrenzte Tests \\
Voisin-Typ-GHR-Detektion & Numerisch & Illustration, kein Kriterium \\
Bedingte Orbit-Rigidität $(\mathrm{AP}'')$ & Schablone (bedingt) & Checkliste \\
Uniformer Anti-Advice-Filter $(\mathrm{AP}''')$ & Methodisches Kriterium & Zulässigkeitsfilter \\
Scale-/Reef$_E$-Provenienz-Gates & Methodischer Guardrail & künftige Matrixchecks \\
\midrule
\multicolumn{3}{l}{\textit{Offen (äquivalent zu Deligne 1982):}} \\
AbsHodge $\stackrel{?}{=}$ Algebraisch & Offen & nicht gelöst \\
Stärkere Positivität jenseits HR & Forschungsrichtung & Folgearbeit \\
\bottomrule
\end{tabular}}
\end{center}


% ===================================================================
% 10. SCHLUSSFOLGERUNG
% ===================================================================
\section{Schlussfolgerung}
\label{sec:conclusion}

Wir haben eine Umformulierung der Hodge-Vermutung durch die Perspektive von Positivität und No-Go-Obstruktionen vorgeschlagen. Das Schlüsselkonzept -- \textit{arithmetische Positivität} -- kombiniert Hodge-Riemann-Positivität, absolute Stabilität unter $\mathrm{Aut}(\mathbb{C})$ und Lefschetz-Kompatibilität in eine einzige Bedingung.

Ein zentrales Ergebnis dieser Arbeit (Bemerkung~\ref{rem:ap_equals_absolute}) ist, dass der arithmetisch positive Kegel mit Delignes absoluten Hodge-Klassen übereinstimmt: $\mathrm{AP}^p(X) = \mathrm{AbsHodge}^p(X)$. Die Arithmetische Positivitätsvermutung ist daher äquivalent zu Delignes Frage ,,Sind alle absoluten Hodge-Klassen algebraisch?'' Die Positivitätsbedingungen (AP1), (AP2b) und (AP3) werden für jede absolute Hodge-Klasse automatisch durch die Hodge-Riemann-Bilinearrelationen garantiert.

Obwohl die Arbeit keine genuin neue Vermutung einführt, bietet sie:

\begin{enumerate}
\item \textbf{Eine neue Perspektive:} Die No-Go-Formulierung verwandelt die Hodge-Vermutung von einem Konstruktionsproblem in eine Unmöglichkeitsaussage.

\item \textbf{Ein diagnostisches Werkzeug:} Das Galois-Hodge-Riemann-Spektrum macht explizit, welche Galois-Konjugation Nichtalgebraizität bezeugen könnte.

\item \textbf{Partielle Funktorialitätstests:} Die Arbeit dokumentiert partielle Funktorialitätstests, keine vollständige Funktorialitätstheorie: Der Rückzug entlang endlicher Morphismen ist an den Vergleich der Amplekegel gebunden, primitive Künneth-Produkte liefern nur Produktpolarisierungs-Vorzeichenformeln, und endlicher Vorwärtsschub bleibt offen.
\end{enumerate}

Die verbleibende Herausforderung ist Delignes Frage: zu beweisen, dass absolute Hodge-Klassen algebraisch sind. Künftige Arbeit sollte stärkere Positivitätsbedingungen suchen -- jenseits der Hodge-Riemann-Relationen -- die die Lücke zwischen absoluten Hodge-Klassen und algebraischen Klassen weiter einengen könnten.

\begin{quote}
\textit{,,Die Hodge-Vermutung verlangt nicht von uns, algebraische Zykel zu bauen. Sie verlangt von uns zu zeigen, dass die Geometrie keinen Raum für etwas anderes hat.''}
\end{quote}


% ===================================================================
% FOERDERINFORMATION
% ===================================================================
\section*{F\"orderinformation}
Diese Forschung erhielt keine externe F\"orderung.


% ===================================================================
% BIBLIOGRAPHIE
% ===================================================================
\begin{thebibliography}{30}

\bibitem{Hodge1950}
W.~V.~D.~Hodge (1950).
The topological invariants of algebraic varieties.
\textit{Proceedings of the International Congress of Mathematicians (Cambridge, MA, 1950)}, vol.~1, pp.\ 181--192. American Mathematical Society.

\bibitem{Voisin2002}
C.~Voisin (2002).
\textit{Hodge Theory and Complex Algebraic Geometry I, II}.
Cambridge University Press.
DOI: 10.1017/CBO9780511615344.

\bibitem{GH1978}
P.~Griffiths, J.~Harris (1978).
\textit{Principles of Algebraic Geometry}.
John Wiley \& Sons, New York.

\bibitem{Deligne1982}
P.~Deligne (1982).
Hodge cycles on abelian varieties (Notes by J.~S.~Milne).
In \textit{Hodge Cycles, Motives, and Shimura Varieties}, Lecture Notes in Mathematics 900, pp.\ 9--100. Springer.
DOI: 10.1007/978-3-540-38955-2\_3.

\bibitem{Andre1996}
Y.~Andr\'e (1996).
Pour une th\'eorie inconditionnelle des motifs.
\textit{Publications Math\'ematiques de l'IH\'ES}, 83, 5--49.
DOI: 10.1007/BF02698643.

\bibitem{Charles2013}
F.~Charles (2013).
The Tate conjecture for K3 surfaces over finite fields.
\textit{Inventiones Mathematicae}, 194(1), 119--145.
DOI: 10.1007/s00222-012-0443-y.

\bibitem{Abdulali1994}
S.~Abdulali (1994).
Algebraic cycles in families of abelian varieties.
\textit{Canadian Journal of Mathematics}, 46(6), 1121--1134.
DOI: 10.4153/CJM-1994-063-0.

\bibitem{Hazama1994}
F.~Hazama (1994).
The generalized Hodge conjecture for stably nondegenerate abelian varieties.
\textit{Compositio Mathematica}, 93(2), 129--137.

\bibitem{Moonen1999}
B.~Moonen (1999).
Notes on Mumford-Tate groups.
Preprint, University of Amsterdam.

\bibitem{Voisin2002b}
C.~Voisin (2002).
A counterexample to the Hodge conjecture extended to K\"ahler varieties.
\textit{International Mathematics Research Notices}, 2002(20), 1057--1075.
DOI: 10.1155/S1073792802111135.

\bibitem{CattaniDeligneKaplan1995}
E.~Cattani, P.~Deligne \& S.~Kaplan (1995).
On the locus of Hodge classes.
\textit{Journal of the American Mathematical Society}, 8(2), 483--506.
DOI: 10.2307/2152824.

\bibitem{Grothendieck1969}
A.~Grothendieck (1969).
Standard conjectures on algebraic cycles.
In \textit{Algebraic Geometry (Bombay, 1968)}, pp.\ 193--199. Oxford University Press.

\bibitem{GrossZagier1986}
B.~Gross \& D.~Zagier (1986).
Heegner points and derivatives of $L$-series.
\textit{Inventiones Mathematicae}, 84, 225--320.
DOI: 10.1007/BF01388809.

\bibitem{Deligne1974}
P.~Deligne (1974).
La conjecture de Weil: I.
\textit{Publications Math\'ematiques de l'IH\'ES}, 43, 273--307.
DOI: 10.1007/BF02684373.

\bibitem{FST-Unified2026}
L.~Geiger (2026).
Functional Stability Theory---Physical Instantiation via the Renormalized Free Energy Principle.
Zenodo preprint, Version 1.8.
\href{https://doi.org/10.5281/zenodo.19036190}{doi:10.5281/zenodo.19036190}.

\bibitem{GeigerRH2026c}
L.~Geiger (2026).
From Landscape to Atlas: Multi-Route Cartography of an Ongoing Expedition Toward the Riemann Hypothesis.
Zenodo preprint, Version 3.1.
\href{https://doi.org/10.5281/zenodo.19035640}{doi:10.5281/zenodo.19035640}.

\bibitem{Mazur1972}
B.~Mazur (1972).
Frobenius and the Hodge filtration.
\textit{Bulletin of the American Mathematical Society}, 78(5), 653--667.
DOI: 10.1090/S0002-9904-1972-12976-8.

\bibitem{Moriwaki2000}
A.~Moriwaki (2000).
Arithmetic height functions over finitely generated fields.
\textit{Inventiones Mathematicae}, 140(1), 101--142.
DOI: 10.1007/s002220050358.

\bibitem{DeligneMilne1982}
P.~Deligne \& J.~S.~Milne (1982).
Tannakian categories.
In \textit{Hodge Cycles, Motives, and Shimura Varieties}, Lecture Notes in Mathematics 900, pp.\ 101--228. Springer.
DOI: 10.1007/978-3-540-38955-2\_4.

\bibitem{BhattScholze2022}
B.~Bhatt \& P.~Scholze (2022).
Prisms and prismatic cohomology.
\textit{Annals of Mathematics}, 196(3), 1135--1275.
DOI: 10.4007/annals.2022.196.3.5.

\bibitem{Bridgeland2007}
T.~Bridgeland (2007).
Stability conditions on triangulated categories.
\textit{Annals of Mathematics}, 166(2), 317--345.
DOI: 10.4007/annals.2007.166.317.

\bibitem{BBT2023}
B.~Bakker, B.~Brunebarbe \& J.~Tsimerman (2023).
o-minimal GAGA and a conjecture of Griffiths.
\textit{Inventiones Mathematicae}, 232(1), 163--228.
DOI: 10.1007/s00222-022-01166-1.

\bibitem{Nie2021} T.~Nie, \textit{Absolute Hodge Cycles in Prismatic Cohomology}, PhD Thesis, Harvard University, 2021.

\bibitem{BinyaminiNovikov2017}
G.~Binyamini \& D.~Novikov (2017).
Wilkie's conjecture for restricted elementary functions.
\textit{Annals of Mathematics}, 186(1), 237--275.
DOI: 10.4007/annals.2017.186.1.6.

\bibitem{BhattMorrowScholze2018}
B.~Bhatt, M.~Morrow \& P.~Scholze (2018).
Integral $p$-adic Hodge theory.
\textit{Publications Math\'ematiques de l'IH\'ES}, 128, 219--397.
DOI: 10.1007/s10240-019-00102-z.

\end{thebibliography}

\makeatletter
\let\test@bbl@sw\false@sw
\auto@bib@empty
\makeatother

\end{document}
